% SPDX-License-Identifier: Apache-2.0
% Technical reference; embedded didactic figures retain their own attribution.
\documentclass[11pt,a4paper,parskip=half,headings=normal,numbers=noenddot]{scrartcl}
\usepackage{fontspec}
\setmainfont{Latin Modern Roman}
\setsansfont{Latin Modern Sans}
\setmonofont{Latin Modern Mono}[BoldFont=lmmonolt10-bold.otf]
\usepackage[ngerman]{babel}
\usepackage{microtype}
\usepackage[a4paper,top=24mm,bottom=25mm,left=25mm,right=25mm]{geometry}
\usepackage{amsmath,amssymb,mathtools}
\usepackage{graphicx,xcolor,booktabs,longtable,tabularx,array}
\usepackage{enumitem,caption,float,placeins}
\usepackage{tikz}
\usetikzlibrary{arrows.meta,positioning,fit,backgrounds,calc}
\usepackage[most]{tcolorbox}
\usepackage{scrlayer-scrpage}
\usepackage{hyperref,bookmark}
\definecolor{ink}{HTML}{193349}
\definecolor{teal}{HTML}{087F8C}
\definecolor{lightteal}{HTML}{EDF7F7}
\definecolor{warm}{HTML}{A55F16}
\definecolor{muted}{HTML}{536575}
\hypersetup{colorlinks=true,linkcolor=ink,urlcolor=teal,citecolor=teal,
 pdftitle={Die Graph-Definition von SkillPilot},
 pdfauthor={SkillPilot / enpasos},
 pdfsubject={Graphsemantik, curriculare Sichten, Lernzustand und implizite Erledigung},
 pdfkeywords={SkillPilot, Skill Graph, Lernziele, Frontier, Assessment, Curriculum}}
\setkomafont{disposition}{\normalfont\sffamily\bfseries\color{ink}}
\setkomafont{pageheadfoot}{\normalfont\sffamily\small\color{muted}}
\clearpairofpagestyles
\ihead{SkillPilot · Graph-Definition}
\ohead{\headmark}
\automark{section}
\cfoot{\pagemark}
\setlength{\headheight}{17pt}
\setcounter{tocdepth}{2}
\setcounter{secnumdepth}{2}
\setlist{nosep,leftmargin=*,topsep=4pt}
\setlist[description]{style=nextline,leftmargin=0pt,labelwidth=0pt,labelindent=0pt}
\captionsetup{font=small,labelfont=bf,skip=7pt}
\renewcommand{\arraystretch}{1.18}
\newcolumntype{P}[1]{>{\raggedright\arraybackslash}p{#1}}
\renewcommand{\tabularxcolumn}[1]{>{\raggedright\arraybackslash}p{#1}}
\setlength{\emergencystretch}{2em}
\newcommand{\code}[1]{\texttt{#1}}
\newcommand{\R}{\mathbb{R}}
\newcommand{\N}{\mathbb{N}}
\newcommand{\Pow}{\mathcal{P}}
\DeclareMathOperator{\Atoms}{Atoms}
\DeclareMathOperator{\Sat}{Sat}
\DeclareMathOperator{\Frontier}{Frontier}
\DeclareMathOperator{\Anc}{Anc}
\DeclareMathOperator{\Desc}{Desc}
\DeclareMathOperator{\ReqAtoms}{RequiredAtoms}
\DeclareMathOperator{\Pre}{Pre}
\DeclareMathOperator{\Missing}{Missing}
\DeclareMathOperator{\MissingAtoms}{MissingAtoms}
\DeclareMathOperator{\Due}{Due}
\DeclareMathOperator{\Verified}{Verified}
\DeclareMathOperator{\Consumers}{Consumers}
\DeclareMathOperator{\Complete}{Complete}
\DeclareMathOperator{\Id}{Id}
\DeclareMathOperator{\Title}{Title}
\DeclareMathOperator{\Weight}{Weight}
\DeclareMathOperator{\ShortKey}{ShortKey}
\DeclareMathOperator{\PhaseCompat}{PhaseCompat}
\DeclareMathOperator{\Applicability}{Applicability}
\DeclareMathOperator{\Work}{Work}
\DeclareMathOperator{\Reserve}{Reserve}
\DeclareMathOperator{\Budget}{Budget}
\newcommand{\Discharged}{\mathrm{Discharged}}
\newcommand{\Active}{\mathrm{Active}}
\newcommand{\Score}{\mathit{score}}
\newcommand{\Passed}{\mathit{passed}}
\newcommand{\MaxPoints}{\mathit{maxPoints}}
\newcommand{\PassingPoints}{\mathit{passingPoints}}
\newcommand{\Reff}{R_{\mathrm{eff}}}
\newcommand{\srcref}[2]{\href{https://github.com/enpasos/skillpilot/blob/main/#1}{#2}}
\newtcolorbox{principle}[1]{breakable,colback=lightteal,colframe=teal,
 boxrule=.5pt,arc=1mm,left=3mm,right=3mm,top=2mm,bottom=2mm,
 title={#1},fonttitle=\sffamily\bfseries}
\newcommand{\sourcefig}[4][.48\textheight]{%
 \begin{figure}[H]\centering
 \includegraphics[width=\linewidth,height=#1,keepaspectratio]{assets/graph-definition/#2}
\caption[\csname title#4\endcsname]{#3}\label{#4}\end{figure}}
\expandafter\def\csname titlefig:overview\endcsname{Ziele und zwei Kantenarten}
\expandafter\def\csname titlefig:contains\endcsname{Direkte und indirekte Zugehörigkeit}
\expandafter\def\csname titlefig:requires\endcsname{Direkte Voraussetzungen}
\expandafter\def\csname titlefig:inheritance\endcsname{Vom Cluster geerbte Voraussetzungen}
\expandafter\def\csname titlefig:cycles\endcsname{Effektive Zyklen durch Vererbung}
\expandafter\def\csname titlefig:frontier\endcsname{Die nächsten verfügbaren Ziele}
\expandafter\def\csname titlefig:filters\endcsname{Normalmodus und strikter Modus}
\tikzset{goal/.style={draw=teal,fill=lightteal,rounded corners=2pt,
 align=center,minimum height=10mm,font=\small,inner sep=5pt},
 req/.style={-{Stealth[length=2mm]},thick,draw=teal},
 grouping/.style={-{Stealth[length=2mm]},draw=muted},
 info/.style={align=center,font=\small,text=muted}}

\begin{document}
\begin{titlepage}
\vspace*{12mm}
{\sffamily\large\color{teal} SKILLPILOT · TECHNISCHE REFERENZ\par}
\vspace{14mm}
{\sffamily\bfseries\fontsize{32}{37}\selectfont\color{ink}
Die Graph-Definition\par von SkillPilot\par}
\vspace{9mm}
{\Large Struktur, Mathematik, curriculare Sichten\par und implizite Erledigung\par}
\vspace{14mm}
\begin{center}\begin{tikzpicture}[node distance=14mm]
\node[goal,minimum width=39mm] (c) {Fachlicher Cluster};
\node[goal,below left=14mm and 3mm of c,minimum width=32mm] (a) {Grundlage $a$};
\node[goal,below right=14mm and 3mm of c,minimum width=32mm] (b) {Anwendung $b$};
\draw[grouping] (c) -- node[left,font=\scriptsize] {contains} (a);
\draw[grouping] (c) -- node[right,font=\scriptsize] {contains} (b);
\draw[req] (b) -- node[above,font=\scriptsize] {requires} (a);
\coordinate (basecenter) at ($(a.south)!0.5!(b.south)$);
\node[info,anchor=north,text width=123mm] at ($(basecenter)+(0,-8mm)$)
 {Gruppierung, Lernvoraussetzung und Abschlussnachweis\\sind unterschiedliche Aussagen.};
\end{tikzpicture}\end{center}
\vfill
\begin{principle}{Geltung dieser Fassung}
Deutsche Gesamtdarstellung der technischen Modellverträge, Stand
\textbf{31. Oktober 2026}. Sie verbindet die normative Graph-Definition mit
ihren angrenzenden Konzepten. Optionale Planungsprofile sind als solche
gekennzeichnet. Die allgemeine Bedarfsrechnung, der direkte Prüfungszugang
und die inklusive Plananrechnung sind lokal implementiert. Ihre Beschreibung
ist kein Nachweis eines Deployments oder einer tatsächlichen Coach-Abnahme.
\end{principle}
\vspace{6mm}
{\small SkillPilot · enpasos -- Enterprise Patterns \& Solutions GmbH\par
Technischer Text und Build-Werkzeuge: Apache-2.0.\par
Didaktische Abbildungen: projektbezogene Inhaltslizenz CC BY 4.0; siehe Anhang.}
\end{titlepage}

\section*{Leserführung und Geltungsbereich}
\markboth{Leserführung}{Leserführung}
\addcontentsline{toc}{section}{Leserführung und Geltungsbereich}
Ein SkillPilot-Graph beschreibt, \emph{was} gelernt werden kann und welche
Voraussetzungen dabei gelten. Curriculare Sichten bestimmen, welche dieser
stabilen Ziele für eine Person gelten. Der Lernzustand hält für jedes atomare
Ziel fest, ob es geschafft ist. Planung verteilt Erwerb und Übung auf verfügbare Zeit.
Die automatische Bedarfsrechnung bestimmt aus Mastery und direkten Anforderungen,
ob ein genau benannter Lernumfang abgeschlossen ist. Diese Aussagen dürfen sich aufeinander beziehen,
aber nicht stillschweigend gegenseitig ersetzen.

Diese Referenz ist für die deutschsprachige gemeinsame Prüfung und
Weiterentwicklung des Modells bestimmt. Sie behandelt die Graph-Definition
im weiteren Sinne: Identität, Kanten, Vererbung, Gültigkeit, Verfügbarkeit,
Projektion, Zieltypen, Fortschritt, Zeitplanung, Abschluss, Quellen und
Qualitätssicherung. Fachliche Lehrplaninhalte und konkrete Aufgabenbestände
werden nicht vollständig abgedruckt. Die Beispiele sind Veranschaulichungen,
keine als Curriculum freigegebenen Teilgraphen.

\begin{center}\begin{tabularx}{\linewidth}{@{}P{35mm}X@{}}
\toprule
Status & Bedeutung in dieser Referenz \\
\midrule
Normativer Kern & Verbindliche mathematische Graphsemantik aus
\code{graph-definition.md}; zusätzliche Profile können strengere Regeln setzen.\\
Angrenzender Vertrag & Semantik von Projektion, Zieltypen, Quellen oder
Lernablauf; Implementierungsstand ist jeweils gesondert zu prüfen.\\
Optionales Profil & Curriculum-Level mit Struktur- und Zeitprüfung;
konzeptionell definiert, kein behaupteter neuer produktiver Validator.\\
Aktuelle Umsetzung & Automatische Bedarfsberechnung und direkter Prüfungszugang;
Deployment und Coach-Abnahme bleiben gesonderte Nachweise.\\
\bottomrule
\end{tabularx}\end{center}

Die fachlichen Verträge im Repository bleiben die maßgeblichen
Einzelquellen. Der Anhang ordnet alle Kapitel diesen Quellen zu und nennt
beim Erstellen überprüfte Implementierungsabweichungen. Mathematische
Beispiele, Bildunterschriften und Plausibilitätsprüfungen ersetzen weder
fachliche Aufgabenreviews noch Untersuchungen mit Lernenden.

\clearpage
\ohead{Inhalt und Abbildungen}
\tableofcontents
\bigskip
\listoffigures
\clearpage
\ohead{Formales Graphmodell}
% SPDX-License-Identifier: Apache-2.0
\section{Ziele, Identität und Beziehungen}
\label{sec:core}
\subsection{Das Grundobjekt}
Der normative Kern ist ein endlicher gerichteter Graph
\[
 \mathcal G=(G,C,R),\qquad C,R\subseteq G\times G.
\]
$G$ ist die Menge der Lernziele. $C$ steht für \code{contains},
$R$ für \code{requires}. Es gibt genau diese beiden \emph{ausdrücklich
verfassten} Kantenarten im Zielgraphen. Abgeleitete Relationen, curriculare
Platzierungen, Quellenzuordnungen und Nachweise sind andere Modellschichten.

Ein Paar $(x,y)$ ist gerichtet. Für jede Relation $X$ bedeutet
\[
 X^{-1}=\{(y,x)\mid(x,y)\in X\}.
\]
$(x,y)\in X^+$ gilt genau dann, wenn ein nichtleerer gerichteter $X$-Pfad
von $x$ nach $y$ führt. $X^+$ ist die transitive Hülle;
Pfade der Länge null zählen nicht.
Ein gerichteter Graph ist azyklisch genau dann, wenn
$\forall g\in G:(g,g)\notin X^+$. Die Vereinigung $C\cup R$ ist
\emph{nicht} die Relation, an der pauschal alle DAG-Bedingungen geprüft werden:
Beide Kantenarten haben verschiedene Bedeutungen.

\begin{principle}{Die Pfeilrichtung ist Teil der Definition}
$p\xrightarrow{\texttt{contains}}c$: Der Cluster $p$ enthält das Ziel $c$.
\quad
$g\xrightarrow{\texttt{requires}}p$: Das Ziel $g$ benötigt $p$.
Eine Lernablaufdarstellung zeigt dagegen $R^{-1}$: von der Voraussetzung
zur darauf aufbauenden Kompetenz. Diese Umkehrung muss beschriftet sein.
\end{principle}

\sourcefig[.42\textheight]{01-skill-graph-overview.png}
{Zwei Beziehungen in einem Graphen. Grau gruppiert, Blau zeigt auf
Voraussetzungen. \emph{Velocity} benötigt \emph{Vectors};
\emph{Newton's Second Law} benötigt sowohl \emph{Forces} als auch
\emph{Acceleration}. Die Farben der Knoten sind hier keine Lernzustände.}
{fig:overview}

\subsection{Attribute und stabile Identität}
Die normativen Basisattribute sind totale Abbildungen
\[
 \Id:G\to\mathrm{UUID},\qquad \Title:G\to\Sigma^*,\qquad
 \Weight:G\to\R_{>0}.
\]
$\Id$ ist injektiv: Verschiedene Ziele haben verschiedene IDs.
Die UUID bleibt über Platzierungen, Sprachfassungen und Ansichten hinweg
die kanonische Identität. Kurze Buchstaben in diesem Dokument sind nur
mathematische Bezeichner, keine Aussage zum serialisierten ID-Format.

Ein optionaler $\ShortKey:G\rightharpoonup\Sigma^*$ ist innerhalb einer
logischen Landschaft eindeutig. Mehrere Sprachdateien desselben Ziels
dürfen dasselbe Paar aus ID und ShortKey wiederholen. Die Teilabbildung
$\PhaseCompat:G\rightharpoonup P_{\mathrm{compat}}$ beschreibt dagegen
Kompatibilitäts- oder Anzeigemetadaten, häufig im Feld \code{phase}.
Sie definiert weder kanonische Identität noch die fachliche Lernreihenfolge.
Weitere gebräuchliche Felder ordnet Tabelle~\ref{tab:attributes} ein.

\begin{table}[H]\centering
\begin{tabularx}{\linewidth}{@{}P{29mm}X@{}}
\toprule
Feld & Fachliche Bedeutung außerhalb des minimalen Kerns \\
\midrule
\code{description} & Verständlicher, überprüfbarer Zieltext; typischerweise
\glqq Die lernende Person kann \ldots\grqq. Orientierung hat eine eigene Semantik.\\
\code{area}, \code{tags} & Fachliche Einordnung und Metadaten; ein Tag ersetzt
keine autoritative Entscheidung über Zieltyp oder Geltungsbereich.\\
\code{level} & Schwierigkeit beziehungsweise Anforderungsbereich;
kein curriculares Erwerbslevel und kein Personalisierungslevel.\\
\code{core} & Kernziel versus Erweiterung im jeweiligen Vertrag;
eine Erweiterung kann im gewählten Profil trotzdem verbindliches Ziel sein.\\
\code{weight} & Anteil bei Fortschrittsaggregation; keine Zeitdauer,
keine fachliche Nachweisstärke und keine Bestehenswahrscheinlichkeit.\\
\code{examples} & Referenzen auf Beispiele oder Aufgaben;
Vorhandensein beweist keine Aufgabenqualität oder Kompetenzbeherrschung.\\
\code{resourceLinks} & Bestehendes Referenzfeld unter anderem für eigene
Visualisierungen, offizielle Quellen und Altlinks. Neue optionale externe
Lernmaterialien gehören nach aktuellem Vertrag in getrennte Inhaltspakete.\\
\bottomrule
\end{tabularx}
\caption{Ergänzende Zielattribute und ihre begrenzte Bedeutung.}
\label{tab:attributes}
\end{table}

\subsection{Atomare Ziele und Cluster}
Die \emph{strukturelle} Klassifikation folgt allein aus $C$:
\[
 A=\{g\in G\mid\nexists c\in G:(g,c)\in C\},\qquad K=G\setminus A.
\]
Ein Blatt ist atomar, jeder andere Knoten ein Cluster. Ein Ziel kann mehrere
Eltern haben: Der kanonische \code{contains}-Graph darf eine Polyhierarchie
sein. Explizite Metadaten zur Klassifikation müssen dazu passen.

Strukturelle Atomarität bedeutet nicht automatisch \emph{semantische}
Atomarität. Ein Blatt mit fünf unabhängigen Teilkompetenzen bleibt formal
ein Blatt, kann jedoch die Qualitätsprüfung verfehlen. Orientierung,
Gedächtnisziele und Prüfungen sind besondere pädagogische Arten; sie sind
keine zusätzlichen Strukturtypen und keine Programmabschnitte.

\section{Containment, Voraussetzungen und Vererbung}
\subsection{Hierarchie und atomare Basis}
Für den endlichen azyklischen Containment-Graphen gelten
\begin{align}
 \Anc(g)&=\{a\in G\mid(a,g)\in C^+\},\\
 \Desc(g)&=\{d\in G\mid(g,d)\in C^+\},\\
 \Atoms(g)&=\begin{cases}\{g\},&g\in A,\\\Desc(g)\cap A,&g\in K.\end{cases}
\end{align}
Die Kurzformen $\Anc$ und $\Desc$ entsprechen den Bezeichnern
\code{Ancestors} und \code{Descendants} der Quellspezifikation und ihrer Bilder.
Jeder Cluster besitzt mindestens ein atomares Blatt: In einem endlichen
azyklischen Graphen endet jeder hinreichend lange Pfad. Deshalb ist
$\Atoms(g)$ für jedes Ziel nichtleer. Ein Knoten ohne Kinder ist atomar,
kein leerer Cluster. Bei mehreren Pfaden zählt ein Atom in der Menge nur einmal.

\sourcefig{02-contains-relation.png}
{Direkte und indirekte Zugehörigkeit. Der türkis hervorgehobene Pfad besteht
aus zwei vorhandenen \code{contains}-Kanten; er ist keine dritte Kantenart.
\emph{Mechanics} ist deshalb auch Vorfahr von \emph{Velocity}.}
{fig:contains}
\FloatBarrier

\subsection{Direkte Anforderungen sind konjunktiv}
$(g,p)\in R$ bedeutet, dass $g$ die Voraussetzung $p$ hat. Mehrere
ausgehende Kanten bedeuten \emph{alle} Voraussetzungen, keine Auswahl.
Alternativen brauchen einen eigenen Auswahlvertrag außerhalb dieser Relation.
Die Relation $R$ muss azyklisch sein.

Als kanonisches Ziel der Modellierung wird $R\subseteq A\times A$
bevorzugt. So liegt fachliche Sequenzierung auf den konkreten Kompetenzen.
Eine am Cluster verfasste Voraussetzung ist stärker als eine summarische
Darstellung: Sie gilt durch Vererbung für alle Nachfahren.

\sourcefig{03-direct-requires.png}
{Direkte Voraussetzungen. Auch der gebogene Pfeil ist eine direkte Kante.
Die Kette von Beschleunigung über Geschwindigkeit zu Vektoren ergibt
transitive Erreichbarkeit, keine Containment-Vererbung.}
{fig:requires}
\FloatBarrier

\subsection{Effektive Anforderungen}
Die effektive Relation entsteht aus eigenen und von Vorfahren geerbten
direkten Anforderungen:
\begin{equation}
 (g,p)\in\Reff\iff
 (g,p)\in R\ \lor\ \exists a\in\Anc(g):(a,p)\in R.
 \label{eq:effective}
\end{equation}
Dabei sind
\[
 R_{\mathrm{inherited}}=\Reff\setminus R,\qquad
 \Reff=R\cup R_{\mathrm{inherited}},\qquad
 \Pre_{\mathrm{eff}}(g)=\{p\mid(g,p)\in\Reff\}.
\]
Mehrere Elternpfade addieren Anforderungen als Mengenvereinigung, ohne
Mehrfachzählung. \emph{Vererbung} in \eqref{eq:effective} und
\emph{transitive Hülle} $\Reff^+$ sind verschiedene Operationen.
Die transitive Hülle beschreibt statische Pfaderreichbarkeit im Graphen.
Laufzeitverfügbarkeit prüft dagegen direkte effektive Anforderungen;
Cluster werden über ihre atomare Basis erfüllt. Mastery einer Voraussetzung
beendet deren eigenen Bedarf, statt weitere Pfadglieder nachzufordern.

Sind alle direkten Anforderungen atomar verfasst, gilt $\Reff=R$:
Containment-Vorfahren sind Cluster und tragen dann keine ausgehenden
Anforderungen. Transitive Ketten bleiben als Graphstruktur relevant.

\sourcefig{04-effective-requires.png}
{Eine Cluster-Anforderung wird vererbt. \emph{Vectors} liegt in diesem
Beispiel außerhalb von \emph{Kinematics}. Würde es auch darin enthalten sein,
erbte es die Anforderung an sich selbst. Durchgezogen heißt direkt verfasst,
gestrichelt heißt geerbt.}
{fig:inheritance}
\FloatBarrier

\section{Gültigkeit, Minimalität und Lernbarkeit}
\label{sec:validity}
\subsection{Verbindliche Bedingungen}
Ein vollständiger Graph ist unter dem normativen Kern genau dann gültig,
wenn die Attribut- und Identitätsbedingungen sowie die folgenden Bedingungen
erfüllt sind:
\begin{align}
 &\forall g\in G:\ (g,g)\notin C^+,\quad (g,g)\notin R^+,
                  \quad (g,g)\notin\Reff^+,\label{eq:acyclic}\\
 &(g,p)\in R\Longrightarrow
      \nexists a\in\Anc(g):(a,p)\in R.\label{eq:localmin}
\end{align}
Die frühere Pflicht zur \emph{transitiven Minimalität} entfällt mit der
Abschlussentscheidung vom 31. Oktober 2026. In $C\to B\to A$ beendet Mastery
von B den Bedarf an A; eine zusätzliche Kante $C\to A$ bleibt dagegen eine
eigene Anforderung. Eine Kante darf deshalb nicht allein aufgrund unveränderter
Erreichbarkeit entfernt werden. Identisch geerbte direkte Paare bleiben lokal
redundant; Zyklen bleiben unzulässig.

Die DAG-Bedingung für $R$ folgt bereits aus derjenigen für $\Reff$, weil
$R\subseteq\Reff$. Sie bleibt als verständliche eigenständige
Autorenprüfung aufgeführt. Unabhängig davon bleibt die Azyklizität von $C$
erforderlich. Identitäten und Referenzen müssen auflösbar sein; unbekannte
Knoten sind keine erfüllten Voraussetzungen.

\sourcefig{05-validity-and-cycles.png}
{Direkte Azyklizität genügt nicht. Aus $A$ enthält $B$ und $A$ benötigt $X$
folgt die geerbte Kante $B\to X$. Kommt $X\to B$ hinzu, entsteht der
effektive Zyklus $B\to X\to B$.}
{fig:cycles}
\FloatBarrier

\subsection{Zusatzprofile und Gegenbeispiele}
Ein Baumprofil kann höchstens einen Elternknoten fordern:
\[
 \forall c\in G:\quad
 \bigl|\{p\in G\mid(p,c)\in C\}\bigr|\leq1.
\]
Dies ist keine Grundbedingung der Polyhierarchie. Ebenso sind Anforderungen
an eigene Vorfahren unerwünscht und in bestimmten Validatorprofilen
verboten. Anforderungen an eigene Nachfahren erzeugen durch Vererbung sogar
eine effektive Selbstabhängigkeit und scheitern bereits am DAG-Vertrag.

\begin{principle}{Ein DAG ist noch kein Beweis der Lernbarkeit}
Seien $G=\{k,a\}$, $C=\{(k,a)\}$ und $R=\{(a,k)\}$.
Die direkten und effektiven Relationen sind azyklisch und minimal.
Aber $\Atoms(k)=\{a\}$: Um $a$ zu beginnen, müsste $a$ schon erfüllt sein.
Die globale Frontier aus leerem Lernzustand ist leer.
Zusätzliche Verbote von Anforderungen an Vorfahren oder ein echter
Abschlusszeuge decken diesen Fehler auf.
\end{principle}

Deshalb müssen drei Aussagen getrennt bleiben: Grundgraph gültig,
gewählte Projektion gültig und Zielmenge aus einem deklarierten Eintrittszustand
vollständig erlernbar. Ein grüner einzelner Validator beweist nicht
automatisch alle drei Aussagen.

\subsection{Didaktische Routen von Orientierung zu Autonomie}
Seien $O\subseteq A$ die Orientierungsanker, $Z\subseteq A$ die
Endziele selbstständiger Leistung und $E_{\mathrm{route}}\subseteq A$
explizit ausgenommene Unterstützungsziele. Die direkte atomare Relation sei
% O and Z avoid collisions with mastery and target sets.
$R_A=R\cap(A\times A)$. Ohne Ausnahmeprofil gilt
$E_{\mathrm{route}}=\varnothing$; Ausnahmen benötigen einen dokumentierten,
maschinenlesbaren Auswahlvertrag. Ein Atom $a\in A\setminus E_{\mathrm{route}}$
ist routenabgedeckt, wenn
\begin{equation}
 \begin{split}
 \exists o\in O,\ z\in Z:\quad
 &(a=o\ \lor\ (o,a)\in(R_A^{-1})^+)\\[-1mm]
 {}\land\;&(a=z\ \lor\ (a,z)\in(R_A^{-1})^+).
 \end{split}
\end{equation}
Die Mengen $O$ und $Z$ heißen in Abschnitt 8.5 der englischen
Quellspezifikation $M$ und $T$. Die abweichenden Buchstaben vermeiden hier
Verwechslungen mit $M_A$ sowie $T_\pi$ und $T_i$.
Die Beschränkung auf $R_A$ verhindert, dass ein Cluster als versteckter
Zwischenknoten einen atomaren Routennachweis liefert. Diese Formel ist
ein Strukturkriterium, kein Beweis didaktischer Angemessenheit.

Konkrete Routenprofile schränken zusätzlich die Schulstufe, erlaubte
semantische Zielarten und die Reihenfolge ein: Orientierung am Einstieg,
fachliches Lernen, bei Bedarf Gedächtnisunterstützung und eine selbstständige
Leistung am Ende. Gedächtnisziele werden dadurch nicht überall Pflicht.
Die Profilregeln \code{CQR-101}, \code{CQR-102}, \code{CQR-103} und
\code{GVR-012}/\code{GVR-013} haben jeweils einen benannten Prüfumfang.
Sie sind nicht als pauschale Zusicherung für jeden Graphen zu lesen.

Eine Übersicht zwischen Clustern kann aus atomaren Routen abgeleitet werden.
Beispielsweise bedeutet eine existenzielle Zusammenfassung $k_2\to k_1$:
\[
 \exists a\in\Atoms(k_2),\ b\in\Atoms(k_1):(a,b)\in R_A^+.
\]
Eine prozentuale Abdeckung braucht darüber hinaus einen erklärten Nenner.
Solche Übersichten erzeugen keine neuen direkten Cluster-Anforderungen.

\section{Lernzustand, Erfüllung und Fortschritt}
\label{sec:state}
\subsection{Erfüllung statt Rückschluss auf Voraussetzungen}
Für die mathematische Progression sei $M_A\subseteq A$ die Menge der
nach dem jeweiligen Zieltyp als erfüllt geltenden atomaren Ziele.
Jedes atomare Lernziel hat genau einen binären Ergebniszustand:
geschafft ($a\in M_A$) oder noch nicht geschafft ($a\notin M_A$).
Teilweise Bearbeitung ist kein dritter Mastery-Zustand.
Die Bezeichnung \glqq mastered set\grqq{} der Spezifikation ist eine
Abstraktion: Bei Orientierung steht Mitgliedschaft für Abschluss, nicht
für fachliche Beherrschung. Bei Gedächtnis- und Prüfungszielen gelten die
jeweiligen besonderen Nachweisregeln.
\begin{equation}
 \Sat(g,M_A)\iff
 \Atoms(g)\ne\varnothing\ \land\ \Atoms(g)\subseteq M_A.
 \label{eq:sat}
\end{equation}
Eine bereits erfolgreiche Anwendung impliziert keine gespeicherte
Beherrschung ihrer Voraussetzungen. $M_A$ wird nicht automatisch
voraussetzungsabgeschlossen. Auch ein als erfüllt angezeigter Cluster
berechtigt nicht zum pauschalen Schreiben von Nachweisen seiner Kinder.

\subsection{Globale Frontier}
Die Menge der nächsten verfügbaren atomaren Ziele lautet
\begin{equation}
 \Frontier(M_A)=\left\{g\in A\setminus M_A\;\middle|\;
 \forall p\in G:\ (g,p)\in\Reff\Rightarrow\Sat(p,M_A)\right\}.
 \label{eq:frontier}
\end{equation}
Sie fordert die direkten effektiven Voraussetzungen. An gemeisterten
Voraussetzungen endet die weitere Bedarfsverfolgung. Ein teilweise
bearbeitetes Ziel ist nicht schon deshalb erfüllt. Ein Cluster als
Voraussetzung wird über seine atomare Basis geprüft.

Für Navigation kann getrennt eine Clusterverfügbarkeit definiert werden:
\begin{equation}
 \begin{split}
 \Frontier_K(M_A)=\{k\in K\mid&\ \Atoms(k)\ne\varnothing,
 \ \neg\Sat(k,M_A),\\
 &\forall p\in G:\ (k,p)\in\Reff\Rightarrow\Sat(p,M_A)\}.
 \end{split}
\end{equation}
Sie ersetzt die atomare Lernfrontier nicht. Eine Oberfläche kann darüber
hinaus verbindliche Zugangs- und Auswahlregeln anwenden, etwa den
Orientierungseinstieg. Die mathematische Frontier ist deshalb nicht
automatisch identisch mit jedem angezeigten UI-Menü.

\sourcefig{06-available-next-goals.png}
{Mit $M_A=\{\text{Vectors},\text{Forces}\}$ ist nur \emph{Velocity} als
nächstes verfügbar. \emph{Acceleration} wartet auf \emph{Velocity};
\emph{Newton's Second Law} zusätzlich auf \emph{Acceleration}.
Die Cluster sind keine direkt auszuwählenden atomaren Lernschritte.}
{fig:frontier}
\FloatBarrier

\subsection{Aggregation ohne Mehrfachzählung}
Sei $S\subseteq A$ eine fest benannte Menge von Zielatomen,
$w(a)>0$ ihr Fortschrittsgewicht und
$m(a)=\mathbf{1}_{\{a\in M_A\}}\in\{0,1\}$ ihr binärer
Mastery-Indikator. Der gewichtete Fortschritt ist
\begin{equation}
 P(S)=\frac{\sum_{a\in S}w(a)m(a)}{\sum_{a\in S}w(a)}\qquad(S\ne\varnothing).
\end{equation}
Für eine leere Zielmenge gibt es keinen solchen Quotienten; die Anzeige
braucht \glqq keine Ziele im gewählten Umfang\grqq{} statt erfundener 100\,\%.
Für mehrere Cluster wird $S$ als Vereinigung ihrer Zielatome gebildet.
Cluster und Kinder dürfen nicht zusätzlich gemeinsam gezählt werden.
Ein Fortschritt von beispielsweise 50\,\% bezeichnet einen Anteil
geschaffter Ziele im Umfang, keine halbe Beherrschung eines einzelnen Ziels.

Bei der Konvention $w(a)=1$ entspricht das Clustergewicht der Zahl seiner
eindeutigen atomaren Nachfahren. Das ist eine zulässige konkrete
Autorenkonvention, keine Forderung gleicher Gewichte für jede Domäne.
Kern- und Erweiterungsanteile können getrennt ausgewiesen werden; der
verbindliche Personal-Curriculum-Zielumfang entscheidet darüber, was zum
Abschluss gehört. Eine aggregierte Zahl ist weder eine kalibrierte
Wahrscheinlichkeit noch eine Schulnote. Ein Notenmodell bräuchte eine
gesonderte, begründete Abbildung.

\section{Filter, Fokus und globale Diagnose}
\label{sec:filters}
\subsection{Filter als Prädikat und kompilierte Anwendbarkeit}
Ein Filter ist zunächst nur $F:G\to\{0,1\}$ mit
$G_F=\{g\mid F(g)=1\}$. Für dimensionale Anwendbarkeit sei $D$
die Dimensionsmenge und $V_d$ das Vokabular von Dimension $d$. Eine
formal typisierte Menge partieller Zuordnungen ist
\[
 \mathcal V_D=\bigsqcup_{J\subseteq D}\prod_{d\in J}\Pow(V_d),\qquad
 \Applicability:G\rightharpoonup\mathcal V_D.
\]
Die disjunkte Vereinigung hält den jeweiligen Definitionsbereich $J$ fest.
Eine fehlende Dimension ist im allgemeinen Metadatenfilter uneingeschränkt;
eine vorhandene Dimension mit leerer Wertemenge lässt dagegen keinen
konkreten Auswahlwert zu. Nur die Abfrage \code{ALL} ignoriert auch diese
Einschränkung. Die Bedeutung fehlender Metadaten ist keine
Erlaubnis, eine ungeklärte persönliche Schulstufenauswahl als vollständig
aufgelöst zu behandeln.

Für eine aktive Auswahl $Q$ kann die Filterprüfung dimensionsweise erfolgen:
\[
 F_Q(g)=1\iff\bigwedge_{d\in D_{\mathrm{aktiv}}}
 \left\{
 \begin{array}{l}
 Q(d)=\mathrm{ALL},\ \text{oder}\ g\notin\operatorname{dom}(\Applicability),\ \text{oder}\\
 d\notin\operatorname{dom}(\Applicability(g)),\ \text{oder}\\
 Q(d)\in \Applicability(g)(d).
 \end{array}\right.
\]
Die Fälle werden nur ausgewertet, wenn die jeweilige Abbildung definiert
ist. \code{ALL} ist ausschließlich ein Abfragewert und wird nicht als
Geltungsbereich gespeichert. Explizite \code{applicabilityOverrides}
dokumentieren begründete Ausnahmen; die Laufzeit verwendet weiterhin die
kompilierte Anwendbarkeit. Eine Sichtbarkeitsentscheidung beweist noch keine
Quellenabdeckung.

\subsection{Erst global vererben, dann einschränken}
Die induzierten Relationen sind
\[
 C_F=C\cap(G_F\times G_F),\quad R_F=R\cap(G_F\times G_F),\quad
 \Reff|_F=\Reff\cap(G_F\times G_F).
\]
Die effektiven Voraussetzungen werden \emph{vor} der Filterung berechnet.
Sonst könnte eine vom ausgeblendeten Vorfahren geerbte Anforderung
verschwinden, obwohl Ziel und Voraussetzung beide im Fokus liegen.
Die atomare Klassifikation wird ebenfalls nicht aus dem gefilterten Baum
neu erfunden: $A_F=A\cap G_F$. Ein Cluster wird durch das Ausblenden
seiner Kinder nicht zu einem neuen atomaren Ziel.

\subsection{Normalmodus und strikter Modus}
Im Normalmodus (\emph{optimistic}) gilt die relative Basis
\[
 \Atoms_F(g)=\Atoms(g)\cap G_F,\qquad
 \Sat_F(g,M_A)\iff\Atoms_F(g)\subseteq M_A.
\]
Die leere relative Basis blockiert hier nicht. Das bescheinigt keine globale
Kompetenz und verändert keinen gespeicherten Lernzustand. Die Frontier ist
\begin{equation}
 \begin{split}
 \Frontier_{\mathrm{opt}}(M_A,F)=\{g\in A_F\setminus M_A\mid
 \forall p\in G_F:\quad\\
 (g,p)\in(\Reff|_F)\Rightarrow\Sat_F(p,M_A)\}.
 \end{split}
\end{equation}
Im strikten Modus bleiben die Kandidaten im Fokus, die Anforderungen
werden jedoch global geprüft:
\begin{equation}
 \Frontier_{\mathrm{strict}}(M_A,F)=
 \{g\in A_F\setminus M_A\mid
 \forall p\in G:\ (g,p)\in\Reff\Rightarrow\Sat(p,M_A)\}.
\end{equation}
Somit gilt für dieselben Eingaben
\[
 \Frontier_{\mathrm{strict}}(M_A,F)\subseteq
 \Frontier_{\mathrm{opt}}(M_A,F).
\]
Jede eingeschränkte direkte Anforderung ist auch eine globale Anforderung; globale Erfüllung impliziert
die relative Erfüllung. Der umgekehrte Schluss gilt nicht.

Bei $g\to x\to p$ und $G_F=\{g,p\}$ ignoriert der Normalmodus
die außerhalb liegende direkte Voraussetzung x. Der strikte Modus prüft x
global. Ist x gemeistert, verlangt auch er nicht zusätzlich p:
Mastery beendet die weitere Bedarfsverfolgung in beiden Modi.

\sourcefig{07-filter-modes.png}
{Identischer Fokus, unterschiedlicher Prüfbereich. \emph{Vectors} liegt
außerhalb; $M_A=\varnothing$. Normalmodus ignoriert die orange markierten
Anforderungen und erlaubt \emph{Velocity}; strikter Modus lässt keines der
beiden Atome zu. Die Bildverweise 11.3 und 11.4 nennen Abschnitte der
englischen Quellspezifikation.}
{fig:filters}
\FloatBarrier

\subsection{Fehlende Voraussetzungen und fehlende Atome}
Die globale Diagnose trennt unmittelbar benötigte Knoten und deren noch offene Zielatome:
\begin{align}
 \Missing(g,M_A)&=\{p\in G\mid(g,p)\in\Reff\land\neg\Sat(p,M_A)\},\\
 \MissingAtoms(g,M_A)&=
 \left(\bigcup_{p\in \Missing(g,M_A)}\Atoms(p)\right)\setminus M_A.
\end{align}
Für beide Mengen $U$ gelten
$U_{\mathrm{in}}=U\cap G_F$ und $U_{\mathrm{out}}=U\setminus G_F$.
Liegt ein benötigter Cluster im Fokus, sein einziges Atom $b$ aber außerhalb,
kann $\Missing_{\mathrm{out}}=\varnothing$ und dennoch
$\MissingAtoms_{\mathrm{out}}=\{b\}$ gelten. Die Clusterliste allein
verschweigt diese externe Kompetenzlücke.

Die Diagnose folgt der globalen Relation. Sie ist weder die Blockerliste
des Normalmodus noch eine Frontier und erzeugt keine Kanten.

Ein optionaler Prüfbereich $S\subseteq G$ ergibt
\begin{equation}
 \Frontier_{S}(M_A,F)=\{g\in A_F\setminus M_A\mid
 \forall p\in S:\ (g,p)\in\Reff\Rightarrow\Sat(p,M_A)\}.
\end{equation}
$S=G$ ergibt den strikten Modus, $S=G_F$ im Allgemeinen \emph{nicht} den
Normalmodus: Die Erfüllung erforderlicher Cluster verwendet weiterhin ihre
globale atomare Basis statt der auf den Fokus eingeschränkten Basis.
Es werden auch hier direkte effektive Anforderungen geprüft, keine zusätzliche
transitive Hülle. Dieses abgeschwächte strikte Profil ist ein Modellvorschlag,
keine aktivierte dritte Betriebsart.

\noindent\small Quellen dieses Teils:
\srcref{docs/concept/skill-graph/graph-definition.md}{Graph Definition};
\srcref{docs/qa-ci/graph-validation-rules.md}{Graph Validation Rules};
\srcref{AGENTS.md}{Projektkonventionen}.\normalsize

\clearpage
\ohead{Curriculare Sichten und Lernzustand}
% SPDX-License-Identifier: Apache-2.0
% Technical graph documentation; included by the main compendium.

\section{Stabile Ziele, kontextabhängige Struktur}
\label{sec:context-model}

Der Graph beantwortet, \emph{was} eine lernende Person können soll und
welche fachlichen Voraussetzungen dafür benötigt werden. Ein Bildungsgang
beantwortet zusätzlich, \emph{wo} diese Ziele verortet sind. Ein persönlicher
Lernplan bestimmt, \emph{wann} daran gearbeitet werden soll. Diese drei
Fragen brauchen unterschiedliche Daten: Das Umordnen eines Lehrplans darf
weder eine neue Kompetenz erzeugen noch bereits gespeicherte Ergebnisse
löschen.

Das allgemeine Zielmodell unterscheidet vier semantische Ebenen
und eine Darstellungsebene (Tabelle~\ref{tab:layers}). Es beschreibt die Richtung der schrittweisen
Migration; nicht jede hier genannte Datenstruktur ist damit bereits in
allen Landschaften vollständig umgesetzt.

\begin{table}[htbp]
\centering\small
\begin{tabularx}{\linewidth}{@{}P{27mm}XX@{}}
\toprule
Ebene & Bedeutung & Beispiel und Grenze \\
\midrule
Ziele & Konkrete Kompetenzen und fachliche Bündel; hier liegen
\code{requires} und \code{contains}. & Eine lineare Gleichung lösen;
Grundlagen linearer Funktionen. \\
Programme & Jahre, Semester, Module, Phasen und andere strukturelle
Einheiten. & Jahrgang 8 ist ein organisatorischer Abschnitt, keine
Kompetenz. \\
Platzierungen & Kontextuelle Zuordnung eines Ziels zu einer
Programmeinheit. & Dasselbe Ziel wird in einem Bildungsgang in Jahrgang 8,
in einem anderen in Jahrgang 9 erworben. \\
Kompetenzachsen & Fachliche Taxonomien und querschnittliche
Kategorien. & K1 \glqq Mathematisch argumentieren\grqq{} ist eine Kategorie;
\glqq Argumentationslücken erkennen\grqq{} kann ein konkretes Ziel sein. \\
Kompositionssicht & Geprüfter Aufbau des sichtbaren Baums für einen
aufgelösten Bildungsumfang. & Eine Landesansicht verweist auf kanonische
Teilbäume, ohne deren Ziele zu kopieren. \\
\bottomrule
\end{tabularx}
\caption{Fünf Modellierungsaufgaben, ein gemeinsamer Bestand stabiler Ziele.}
\label{tab:layers}
\end{table}

\subsection{Identität trotz unterschiedlicher Einordnung}

Die Identitäts-, Attribut- und Atomaritätsregeln aus
Abschnitt~\ref{sec:core} gelten in allen Schichten unverändert.
Eine neue Platzierung erzeugt deshalb keine neue Kompetenz.
Ein Programmabschnitt vom Typ \code{exam} ist noch kein
Prüfungsziel mit Aufgaben, Bewertungsraster und Lernergebnis.
Und \code{phase: Q1} erzeugt weder eine Elternkante zu einem
Q1-Knoten noch eine Voraussetzung zwischen Q1 und Q2.

Historische Bäume können Jahrgangs- und Phasenknoten als
Kompatibilitätsgerüst enthalten. Im allgemeinen Zielmodell
übernimmt explizite Programmstruktur diese Aufgabe; die
fachlichen Ziele bleiben eigenständig. Der gleiche Inhalt
kann so in mehreren Bildungsgängen anders eingeordnet werden,
ohne seine Ziel-ID oder gespeicherte Ergebnisse zu kopieren.

\subsection{Platzierungen und Projektionen}

\code{programUnits} beschreiben im Zielmodell einen expliziten Baum oder
Wald von Programmeinheiten. \code{goalPlacements} verbinden Ziele mit
diesen Einheiten. Eine \code{primary}-Platzierung liefert den normalen
Programmanchor; \code{secondary} zeigt zusätzliche Relevanz,
\code{assessed} einen Prüfungsbezug. Sekundäre und prüfungsbezogene
Zuordnungen erzeugen keine weiteren sichtbaren Kopien desselben Ziels.

Für einen aufgelösten Kontext gilt: Genau eine passende primäre
Platzierung ergibt einen eindeutigen Anker. Ohne passende Platzierung
darf die Darstellung keinen Elternknoten erfinden. Mehrere passende
primäre Platzierungen benötigen eine ausdrücklich geprüfte Auflösung;
sonst ist die Projektion uneindeutig. Kontextmerkmale wie Bundesland,
Schulform, Stufe, Dauer und Kursprofil ändern dabei nicht die Ziel-ID.

Die normale Inhaltssicht folgt ausschließlich den expliziten
\code{contains}-Kanten. Die Programmsicht folgt Programmeinheiten und
Platzierungen. Die Kompetenzsicht folgt Katalogeinträgen und
Kompetenzreferenzen. Die Reihenfolge der Kinder kommt aus der dafür
deklarierten Datenreihenfolge; eine topologische Sortierung nach
Voraussetzungen ist eine gesondert benannte didaktische Ansicht.

\subsection{Geprüfte Kompositionssichten}

Eine Kompositionssicht beschreibt den sichtbaren oberen Aufbau eines
Curriculums durch Strukturknoten, Beschriftungen, Reihenfolge und
Referenzen auf kanonische Teilbäume. Sie enthält keine neu verfassten
Atomziele und keine neuen \code{requires}-Kanten. Wo eine solche Sicht
für den Kontext maßgeblich ist, darf ein Renderer sie nicht durch
Vermutungen aus Titeln, \code{phase} oder alternativen Platzierungen
überschreiben.

Der aufgelöste Standardbaum muss jedes sichtbare Ziel höchstens einmal
und unter höchstens einem sichtbaren Elternknoten zeigen. Ein globaler
DAG kann mehrere fachliche Elternbezüge besitzen; die sichtbare
Baumprojektion benötigt dafür eine deterministische, geprüfte
Auflösungsregel. Überlappende Teilbaumreferenzen sind ohne einen solchen
Vertrag ungültig. Eine Rollenregel für überlappende Referenzen ist noch
keine Erlaubnis, Ziele im sichtbaren Baum zu duplizieren.

Die Kompilation soll früh erfolgen: Kontext auflösen, passende Sicht
auswählen, kanonische Referenzen expandieren, Rollen und sichtbare Eltern
prüfen, dann das Ergebnis für die Laufzeit bereitstellen. Damit wird
Strukturvalidierung möglich, bevor eine Benutzeroberfläche daraus einen
Baum zeichnet.

Ein autoritativ als \code{curricularArea} klassifizierter Cluster
darf nicht über eine direkte \code{goalEntry}-Referenz als
undurchsichtiges, atomar wirkendes Laufzeitziel erscheinen. Eine
\code{canonicalSubtree}-Referenz erhält seine geprüften Nachkommen;
alternativ benennt eine explizite Sichtstruktur die beabsichtigten
Atomziele. Diese Darstellungsregel verändert keine Projektionsrolle.

\paragraph{Quellen.}
\srcref{docs/concept/skill-graph/general-goal-system-and-migration.md}{Allgemeines Zielsystem und Migration};
\srcref{docs/concept/skill-graph/view-projection-and-goal-placement.md}{Sichten und Zielplatzierung}.

\section{Persönlicher Umfang, Fokus und sichtbare Ziele}
\label{sec:personal-scope}

\subsection{Vier Personalisierungsstufen}

Die Auswahl \glqq Schule\grqq{}, \glqq Hochschule\grqq{} oder \glqq Sprachen\grqq{} ist zunächst
Katalognavigation. Erst danach beginnen die vier
Personalisierungsstufen. Sie sind von den curricularen Levels zu
unterscheiden, die Erwerbsetappen mit Voraussetzungen und Zeitbudgets
bezeichnen.

\begin{enumerate}
\item \textbf{Basiscurriculum:} ein ausdrücklich auswählbares
Wurzelcurriculum stellt den Zielbestand bereit.
\item \textbf{Personal Curriculum:} der längerfristig festgelegte,
backendseitig verantwortete Bildungsumfang. Dazu gehören beim
Schulcurriculum der Landes- beziehungsweise kanonische Kontext,
gegebenenfalls G8/G9, eine ausdrücklich bestimmte Stufe, die Fächer
und deren Kursprofile.
\item \textbf{Lernfokus und aktives Ziel:} Ein oder mehrere Wurzeln
begrenzen vorübergehend den Arbeitsbereich; genau ein atomarer
Lerngegenstand ist das aktuelle Ziel. Fokus und aktives Ziel sind
veränderlich.
\item \textbf{Mastery:} Das binäre Ergebnis \glqq geschafft\grqq{} oder
\glqq noch nicht geschafft\grqq{} gehört zur stabilen atomaren Ziel-ID.
Die Ergebnisse bleiben bei Änderungen der Sicht oder des
Fokus erhalten. Clusterfortschritt wird daraus abgeleitet.
\end{enumerate}

Ein Kursprofil wie LK ist ein Fachattribut und bestimmt nicht implizit
die Schulstufe. Fehlende Angaben zur Stufe bleiben ungeklärt. Eine
Änderung des Personal Curriculum prüft Fokus und aktives Ziel erneut
gegen die neue Zielmenge; sie darf globale Ergebnisse weder löschen
noch neu vergeben. Die aktuelle Konfiguration des Personal Curriculum
erfolgt über die dafür vorgesehenen SkillPilot-Oberflächen. Der
OpenAI-V1-Coach ändert diese Konfiguration nicht selbständig.

Alle \code{target}-Ziele im Personal Curriculum bleiben beabsichtigte
Lernziele. Ein enger Fokus reduziert den vorübergehenden Arbeitsbereich,
nicht den längerfristigen Gesamtumfang oder dessen Abschlussbedingungen.

\subsection{Zielrolle und Voraussetzungshilfe}

Eine Referenz in einer Kompositionssicht kann ausdrücklich eine von
zwei Rollen tragen:

\begin{description}
\item[\code{target}] Das Ziel gehört zur auswählbaren Lernmenge,
zur Frontier, zum Fortschritt und zum Abschluss dieser Sicht.
\item[\code{prerequisiteOnly}] Das Ziel bleibt für
Voraussetzungsprüfungen mit seiner globalen Identität und seinem
gespeicherten Zustand verfügbar. Es ist in dieser Sicht kein
auswählbares Ziel und zählt nicht zum Fortschrittsnenner.
\end{description}

Fehlt eine Rollenangabe, gilt \code{target}. Teilbaumreferenzen vererben
ihre Rolle an Nachkommen. Bei überlappenden Rollenfestlegungen entscheidet
die Spezifität: Eine direkte \code{goalEntry}-Referenz geht einer
Teilbaumvererbung vor; eine nähere Teilbaumwurzel geht einer breiteren
vor; erst bei gleicher Spezifität gewinnt \code{target}. Diese Regel
löst die Rolle, nicht eine ansonsten unzulässige doppelte Darstellung.

Rollen werden nicht aus Alter, Phase oder Voraussetzungspfaden erraten.
Ein mathematisches Ziel kann in Mathematik ein Lernziel und in Physik
eine unterstützende Voraussetzung sein. Eine gemeinsame Planung
identifiziert es dennoch nur einmal. Seine Physikrolle löscht keine
Mathematikergebnisse und erzeugt auch keinen zweiten Kompetenzerwerb.

Jede direkte kanonische Voraussetzung eines sichtbaren Ziels muss
in der maßgeblichen Level-2-Komposition als \code{target} oder
\code{prerequisiteOnly} auflösbar sein. Fehlt sie, ist die
Konfiguration fehlerhaft und das Ziel bleibt gesperrt. Nur bei
ausgelassenen Anforderungen, die von übergangsweise vorhandenen
Legacy-Clustern geerbt werden, bleibt das besondere
Kompatibilitätsverhalten bestehen. Diese Gültigkeitsprüfung ist
von der späteren Level-3-Fokuswahl getrennt: Eine gültig
aufgelöste Voraussetzung außerhalb des Fokus blockiert den
Normalmodus weiterhin nicht.

\subsection{Anwendbarkeit ist keine Platzierung}

\code{applicability} ist kompilierte Sichtmetadateninformation, etwa
für einen Bundeslandfilter. Sie wird aus überprüfbaren Quellen,
Mappings, Provenienz und begründeten Ausnahmen gewonnen. Sie beantwortet,
in welchem Kontext ein Ziel sichtbar ist. Eine Platzierung beantwortet,
wo es in einer Programmstruktur liegt. Keine der beiden Angaben ist
ein Beweis für individuelle Beherrschung.

Bei landesspezifischen Ergänzungen kann eine ausdrücklich deklarierte
Vererbungsgrenze verhindern, dass eine breite Quellenzuordnung eines
Vorfahren als Quellenbeleg aller neuen Nachkommen ausgegeben wird.
Direkte Quellen und Zuordnungen im Ergänzungszweig bleiben maßgeblich.
Ebenso kann die Sichtbarkeit eines lokalen Übungs- oder Prüfungsendpunkts
unter einem dafür vorgesehenen Vertrag aus seinen Voraussetzungen
abgeleitet werden. Das belegt seine Anwendbarkeit, nicht automatisch
seine Lehrplanabdeckung.

\subsection{Fokuswechsel ohne impliziten Kompetenzerwerb}

Ein Fokus besteht aus der Vereinigung der gewählten Wurzeln und ihrer
sichtbaren Nachkommen, eingeschränkt auf das Personal Curriculum.
Im Normalmodus werden die effektiven Voraussetzungen innerhalb dieses
Fokus geprüft; der explizite strenge Modus prüft global. Eine außerhalb
liegende Voraussetzung wird im Normalmodus weder zum Hindernis noch
dadurch gemeistert. Die formalen Unterschiede der Frontier-Berechnung
werden im mathematischen Teil behandelt.

Eine Fokusausweitung folgt dem eindeutigen sichtbaren Elternpfad. Die
Backendoption muss eine echte Obermenge liefern und mindestens ein
noch offenes Ziel hinzufügen. Bei mehreren Wurzeln ersetzt sie nur den
abgedeckten Zweig und bewahrt die unabhängigen Wurzeln. Der Coach
übernimmt die veröffentlichte Option vollständig und wartet vor der
Änderung auf die Zustimmung der lernenden Person.

Eine automatische Ausweitung wird erst nach tatsächlichem Abschluss
des aktuellen Fokus vorgeschlagen. Eine leere Frontier kann auch einen
blockierten, unvollständigen Fokus bedeuten. Sie beweist keinen
Abschluss. Nach der Ausweitung werden neu enthaltene, unbeherrschte
Voraussetzungen regulär geprüft: Die Beherrschung eines abhängigen Ziels
überträgt keine Mastery rückwärts. Ihre Erforderlichkeit wird aus dem
vollständigen Bedarfskontext neu berechnet; rein implizit erledigte
Erweiterungen begründen keinen neuen Pflichtvorschlag.

\paragraph{Quellen.}
\srcref{docs/concept/levels-of-personalization.md}{Personalisierungsstufen};
\srcref{docs/concept/skill-graph/view-projection-and-goal-placement.md}{Projektionsrollen und Sichtvertrag}.

\section{Zieltypen und die Bedeutung eines Ergebnisses}
\label{sec:goal-types}

Der Graph erhält einen zieltypgerecht aufgelösten Zustand. Er entscheidet
nicht selbst, ob eine Antwort fachlich überzeugend war, ein
Orientierungsgespräch angemessen abgeschlossen wurde oder alle
Karten eines Decks sicher abgerufen wurden. Diese Entscheidungen
gehören zum jeweiligen Ergebnisvertrag.

Das Ergebnis jedes Lernziels bleibt binär: geschafft oder noch nicht
geschafft. Der Zieltyp bestimmt, wann es als geschafft gilt.
Bei einem Inhaltsziel ist dies die fachlich festgestellte Beherrschung;
bei Orientierung der abgeschlossene motivierende Einstieg, kein
fachliches Zertifikat. Bei Memorierung wird das eine Zielergebnis aus
den Kartendaten abgeleitet. Unterschiede der Abschlussbedingungen
erzeugen keine zusätzlichen Mastery-Stufen oder parallelen Zielergebnisse.

\subsection{Verstehen und Anwenden}

Ein gewöhnliches Inhaltsziel beschreibt konzeptuelles oder prozedurales
Verstehen. Die Beurteilung muss die in der Beschreibung verlangten
Aspekte abdecken. Eine gelöste Aufgabe ist ein möglicher Beleg, aber
nicht schon per Definition die Beherrschung des ganzen Ziels.
Mehrschrittiger Transfer in einer Aufgabe kann ausreichend sein; eine
pauschale Mindestzahl von Aufgaben wäre ebenso unbegründet wie das
automatische Anerkennen jeder richtigen Einzelantwort.

Im coachgeführten Ablauf wird eine ausreichend belegte Beherrschung
unmittelbar gespeichert. Erst eine bestätigte Speicherung erlaubt
die Aussage, das Ergebnis sei gesichert. Eine zusätzliche Zustimmung
zum Speichern ist kein fachlicher Nachweis und wird für den regulären
Erfolgsfall nicht als weitere Abschlussbedingung eingeführt.
Der anschließende Themenwechsel bleibt eine andere Entscheidung:
Der Coach bietet das backendseitig gewählte nächste Thema an und wartet
auf erkennbares Einverständnis, bevor neues Material beginnt.

\subsection{Orientierung als eigener Abschlussvertrag}

Ein autoritativ als \code{orientation} klassifiziertes Ziel zeigt
konkrete Möglichkeiten und ehrliche positive Perspektiven des
folgenden Lernstoffs. Es prüft weder Vorwissen noch Fachbegriffe,
Rechenfertigkeiten, Abruf, Transfer oder Prüfungsleistung. Wie bei allen
Lernzielen ist das Ergebnis binär; eine Prozentbewertung der persönlichen
Motivation wäre zudem inhaltlich irreführend.

Wählt die lernende Person eines der angebotenen Interessen, beginnt
dadurch zunächst eine passende motivierende Vertiefung. Der Coach
verknüpft das Interesse mit etwas Konkretem, das die Person verstehen,
erkunden, gestalten oder tun kann, und lädt zu einer einfachen
persönlichen Reaktion, Wahl oder Frage ein. Das bloße Wiederholen
eines Angebotslabels schließt die Orientierung nicht ab.

Nach einer sinnvollen Antwort auf diese Vertiefung oder dem
ausdrücklichen Wunsch, die Orientierung zu verlassen, folgt eine
positive Zusammenfassung mit einem gesonderten Abschlussangebot.
Erst die weitere Zustimmung dazu begründet den Abschlussmarker.
Dieser Ablauf darf weder als Fachprüfung umgedeutet noch auf normale
Inhaltsziele übertragen werden. Orientierungsknoten sind reine
voraussetzungsfreie atomare Einstiege; Prüfungs- und Memory-Metadaten
gehören nicht auf denselben Knoten.

Bei einem echten Neueinstieg ohne bisherigen Fortschritt im projizierten
Umfang sind offene Orientierungsziele die einzigen auswählbaren Frontier-Ziele.
Lernende mit bereits bestehendem Fachfortschritt werden nicht auf diese
Einstiegsregel zurückgesetzt.

\subsection{Memorierung: Übung und verifizierter Abruf}

Ein Memorierziel bündelt die erforderlichen Karten eines Decks.
Die Karten besitzen zwei getrennte Prozesszustände. Die SRS-Übung steuert
Wiederholungsintervalle. Der verifizierte Abruf hält fest, ob eine
Karte in einer kontrollierten Abfrage korrekt beantwortet wurde.
Das Umdrehen einer Karte oder eine Selbsteinschätzung in der Übung
ist kein solcher Nachweis. Aus den Kartendaten folgt genau ein binäres
Mastery-Ergebnis des Lernziels; Übungs- und Abrufprozess sind keine
zusätzlichen gespeicherten Zielergebnisse.

Sei $K_g$ die aufgelöste Menge erforderlicher Karten eines Ziels $g$.
Für einen Zeitpunkt $t$ bezeichnet $\Due(k,t)$ die Fälligkeit und
$\Verified(k)$ einen aktuell gültigen bestandenen Abruf. Der
\emph{dokumentierte Zieltypvertrag} verlangt beide Bedingungen:
\[
  \Complete_{\mathrm{memory,doc}}(g,t)
  \iff
  \bigl(\forall k\in K_g:\neg \Due(k,t)\bigr)
  \land
  \bigl(\forall k\in K_g:\Verified(k)\bigr).
\]
Diese Formel setzt ein vollständig aufgelöstes, gültiges Deck voraus.
Fehlende Kartendaten werden dadurch nicht als leeres, automatisch
abgeschlossenes Deck interpretiert.

\textbf{Abweichung des geprüften Codes.} Der zum Dokumentstand
gelesene \code{LearnerService} und der Frontend-Hook
\code{useSrsMastery} leiten Memorier-Mastery dagegen allein aus dem
verifizierten Abruf aller erforderlichen Karten ab. Für ein nichtleeres,
vollständig aufgelöstes Deck lautet diese Implementierungsregel:
\[
  \Complete_{\mathrm{memory,code}}(g)
  \iff \forall k\in K_g:\Verified(k).
\]
Neue SRS-Fälligkeit entzieht danach nicht automatisch die Mastery.
Ein fehlgeschlagener erneuter Abruf kann sie hingegen zurücknehmen.
Dokumentierte Sollsemantik und aktueller Code sind damit nicht
deckungsgleich; dieses Kompendium trifft keine stillschweigende
Migrationsentscheidung. Die Abweichung ist einmal im Zieltypvertrag
aufzulösen. Die allgemeine Abschlussberechnung verwendet dasselbe binäre
Zielergebnis wie die übrigen Oberflächen; es führt keine eigene
Anerkennung von Memorier-Mastery ein. Fällige Übung bleibt eine
Prozessaussage über die Karten.

Der gegenwärtig dokumentierte SRS-Vertrag verwendet SM-2-Zustand mit
Intervall, Wiederholungszähler, Leichtigkeitsfaktor und nächstem
Wiederholungsdatum. Eine Karte ohne gespeicherten Zustand ist fällig.
Erfolgreiche Übung verlängert nach den festgelegten Regeln das
Intervall; unzureichende Übung setzt die Wiederholungsfolge zurück.
Diese Zeitplanung ersetzt den verifizierten Abruf nicht.

Eine harte Abfrage kann die lernende Person beim aktiven Memorierziel
auch beginnen, wenn die SRS-Warteschlange noch nicht leer ist. Zuerst
wird nur die Frage abgerufen, dann antwortet die Person ohne
Kartenhilfe. Erst anschließend wird die erwartete Antwort geladen
und der Abruf als bestanden oder fehlgeschlagen gespeichert.
Ein bestandener Abruf aktualisiert auch die SRS-Spur als starke
erfolgreiche Wiederholung. Ein Fehlschlag hebt den verifizierten
Status dieser Karte auf und führt sie in die Übung zurück.

\subsection{Prüfungsziele im geltenden Ablauf}

Prüfungsziele verwenden \code{examData} mit Aufgabeninhalt,
Lösungsinhalt und Bewertungsraster. Das Raster enthält maximale
Punkte, Bestehensgrenze und bewertbare Schritte. Für gültig
vergebene Schrittpunkte $p_j$ gilt im aktuellen Vertrag:
\[
  \Score=\min\!\left\{\sum_j p_j,\;\MaxPoints\right\},
  \qquad
  \Passed\iff \Score\ge \PassingPoints.
\]
Die Punktevergabe folgt den einzelnen Rastervorgaben. Die
Gesamtsummenformel allein legitimiert keine frei erfundenen
Punkte, negativen Bewertungen oder Änderungen des Aufgabeninhalts.

Der normale Lernzugang prüft direkte effektive Voraussetzungen einschließlich
der atomaren Erfüllung erforderlicher Cluster. Für Prüfungen gilt zusätzlich
der allgemeine direkte Zugang: Fehlende Vorbereitungsmastery sperrt einen
bewusst gewählten Prüfungsversuch nicht. Aufgabenverfügbarkeit, persönlicher
Geltungsbereich und autorisierte Sitzung bleiben erforderlich. Dieser Zugang
ist lokal implementiert; tatsächliches Deployment und Coach-Abnahme sind
getrennt nachzuweisen.

Im Prüfungsmodus erhält die Person die vollständige, unveränderte Aufgabe
und bearbeitet sie ohne inhaltliche Hilfen. Ihre vollständige Einreichung
wird anhand des Rasters bewertet. Bei Erfolg wird Mastery unmittelbar
gespeichert. Ein Fehlversuch speichert weder Ergebnis noch Punkte und lässt
vorhandene Mastery bestehen. Erst die bestätigte Speicherung erlaubt die
Aussage, der Erfolg sei gesichert.

Nach der Bewertung erhält die Person Ergebnis, Punkte, Begründung und
Lösung; neues Material beginnt nach ihrer weiteren Entscheidung. Die
unbegrenzte Wiederholung derselben Aufgabe bleibt möglich. Bekannte Lösungen
können die Aussagekraft eines Versuchs verringern, erzeugen aber keine zweite
Anerkennungsstufe: Vorhandene Prüfungs-Mastery zählt unmittelbar in der
allgemeinen Abschlussberechnung. Entfallende Arbeit wird aus dem vollständigen
Bedarfsgraphen und der Mastery aller Zielarten abgeleitet, nicht ausschließlich
aus bestandenen Prüfungen.

\paragraph{Quellen.}
\srcref{docs/concept/skill-graph/node-types.md}{Zieltypen und Prüfungsablauf};
\srcref{AGENTS.md}{Projektregeln zu Orientierung, Evidenz und Fortsetzung};
\srcref{backend/src/main/java/com/skillpilot/backend/service/LearnerService.java}{Backend-Ableitung von Memorier-Mastery};
\srcref{app/src/hooks/useSrsMastery.ts}{Frontend-Ableitung von Memorier-Mastery}.

\section{Zeitplanung als ergänzende Sicht}
\label{sec:time-context}

Ein Zeitplan bindet einen erklärten Zielumfang, eine Curriculumfassung,
Kalender, verfügbare Kapazität und Planungsversion. Absolute Termine
gehören nicht auf kanonische Ziele, Cluster oder Platzierungen.
Derselbe Inhalt kann in verschiedenen Kursen zu verschiedenen
Zeiten bearbeitet werden, ohne seine fachliche Identität zu ändern.

Die Zeitplanung ist eine zusätzliche Ausführungsplanung und keine
fünfte Personalisierungsstufe. Sie darf Ziele priorisieren, die im
aktuellen Arbeitsbereich verfügbar sind. Sie erzeugt keine
Voraussetzungskanten, keine neue Frontier und keine Lernergebnisse.
Ein curricularer Level ordnet Erwerbsetappen gröber; ein datierter
Ablaufplan muss darüber hinaus gewährleisten, dass eine Voraussetzung
rechtzeitig vor ihrer geplanten Verwendung erarbeitet wird.

Verfügbare Kapazität berücksichtigt tatsächliche Unterrichts- oder
Lernminuten nach Ferien, Ausfällen und anderen Ausschlüssen.
Aktivitäten, geschützte Reserve und verbindliche Fachbudgets
müssen darin ohne Doppelzählung Platz finden. Die formale
Budgetprüfung und ihre Anforderungen an kalibrierte
Aufwandsschätzungen folgen in Abschnitt~\ref{sec:levels}.

Plan-Soll und tatsächliche Zielerreichung verwenden denselben
benannten Umfang. Gespeicherte Lernergebnisse sind die Quelle
der Zielerreichung; eine zusätzliche manuelle
Unterrichtsabdeckung begründet oder verbirgt keinen Abschluss.
Fehlende Ergebnisse bleiben als nicht verfügbar erkennbar.
Neue Planrevisionen dürfen historische Ziele, Ergebnisse und
frühere Aussagen nicht stillschweigend umdeuten.

Der dokumentierte lokale Planungspilot und die angestrebte
backendseitig maßgebliche Architektur sind getrennt zu betrachten.
Insbesondere behauptet die Leveldefinition noch keine
implementierte fachübergreifende Budgetprüfung.

\paragraph{Quelle.}
\srcref{docs/concept/didactic/curriculum-time-axis-and-pacing.md}{Zeitachse, Kapazität und Planungsgrenzen}.

\section{Herkunft, Materialien und belastbare Qualitätsaussagen}
\label{sec:evidence-context}

\subsection{Drei Arten von Verweisen}

Curriculare Herkunft, eigene Veranschaulichung und optionales
Lernmaterial haben verschiedene Funktionen. Ein offizieller
Quellenverweis begründet die Aufnahme oder Formulierung eines
Lernziels. Ein eigenes Lernzielbild veranschaulicht dessen Kernidee.
Ein Buch, Video oder eine Simulation kann bei der Erarbeitung
helfen. Keiner dieser Verweise ist für sich ein individueller
Kompetenznachweis.

Die verbindliche Inhaltsarchitektur lautet: Materialien
referenzieren das Curriculum; das Curriculum wählt keine
Materialanbieter. Neue optionale externe Lernmaterialien gehören
in getrennte Inhaltspakete, die auf stabile Ziel-IDs verweisen.
Die Auswahl eines Pakets verändert weder Personal Curriculum
noch Voraussetzungen, Mastery oder Abschlussregeln. Unvollständige
Materialabdeckung darf normale Lernarbeit nicht sperren.

Offizielle Quellen, Mappings und strukturierte Provenienz bleiben
am Curriculum nachvollziehbar. Eigene Zielvisualisierungen behalten
ihren gesonderten kanonischen Referenz- und Qualitätsvertrag.
Bestehende externe Links auf kanonischen Zielen sind gegebenenfalls
Migrationsbestand; daraus folgt keine dauerhafte zweite
Autorenarchitektur. Eine reine Linklieferung belegt außerdem
nicht, dass ein Coach den verknüpften Inhalt gelesen hat.

\subsection{Quellenbegründung für Ziele und Beziehungen}

Eine überprüfbare Herkunftskette verbindet das kanonische Ziel mit
einer Quelle, einem konkreten Quellabschnitt, einer nachvollziehbaren
Extraktion und einer geprüften Mappingentscheidung. Diese kann
eins zu eins, teilweise, aggregierend oder auf mehrere Ziele
aufteilend sein. Eine verständliche Begründung erklärt diese Form,
anstatt jede Zuordnung als wortgleiche Lehrplanübernahme darzustellen.

Für \code{requires} und \code{contains} ist ebenfalls zu
unterscheiden, ob eine Beziehung ausdrücklich in der Quelle steht
oder eine fachlich begründete Modellierungsentscheidung ist.
Ein plausibler Zieltext erlaubt keine nachträglich erfundene
Provenienz. Ungeklärte Quellen, veraltete Fassungen und
widersprüchliche Zuordnungen bleiben sichtbare Befunde.

Die menschenlesbare Begründung ist eine Darstellung vorhandener
Evidenz, kein zweites führendes Register. Strukturierte Referenzen
und Fingerprints ermöglichen, dieselbe Begründung in Dokumenten,
Exporten und Laufzeitwerkzeugen konsistent auszugeben. Ein
zusätzlicher MEM/FWU-Abgleich kann die klassische Quellenroute
stützen oder infrage stellen; er ersetzt sie nur unter einem
ausdrücklich geprüften Vertrag für den konkreten Umfang.

\subsection{Algorithmische Gültigkeit und fachliche Qualität}

Die Ebenen der Qualität bauen aufeinander auf, sagen aber nicht
dasselbe aus. Ein Parser erkennt fehlerhaftes JSON; eine
Referenzprüfung unbekannte IDs; ein Zyklustest unerlaubte
Kantenfolgen. Keiner dieser Tests beurteilt allein, ob die
Kompetenz sinnvoll geschnitten, eine Erklärung verständlich
oder eine Aufgabe didaktisch geeignet ist.

Die in Abschnitt~\ref{sec:validity} erläuterten
Gültigkeits- und Routenprofile brauchen jeweils einen benannten
Prüfumfang. Ein lokales Migrationsgate oder ein einzelner
grüner Strukturcheck bestätigt nicht automatisch sämtliche
Anforderungen an Graph, Projektion und Lernbarkeit.

Für vertiefte Curriculum-QS unterscheidet SkillPilot fünf
gleichzeitig aktuelle Nachweisbereiche:
\begin{description}
\item[D -- Beschreibung:] Ist der Zieltext fachlich korrekt und klar?
\item[P -- positive Verständnisnachweise:] Welche beobachtbare
Leistung zeigt das geforderte Verständnis?
\item[A -- Atomarität:] Ist die Kompetenz sinnvoll zugeschnitten?
\item[M -- Memory:] Ist die Entscheidung über notwendige
Memorieranteile begründet und aktuell?
\item[V -- Visualisierung:] Ist die Darstellung fachlich geeignet
und an die geprüften Bild- und Zielstände gebunden?
\end{description}

Der Reifegrad M7 bezeichnet den vollständigen geltenden
maschinellen Abschluss dieser Anforderungen und der
vorausgesetzten unteren Reifegrade für den erklärten Umfang.
Er ist keine menschliche Freigabe, keine menschliche Erprobung
und kein Nachweis von Lernwirksamkeit. Solche Aussagen
benötigen eigene Prüfungen und eigene Evidenz. Auch mehrere
übereinstimmende KI-Reviews ersetzen keine gesondert
vorgeschriebene menschliche Entscheidung.

Eine Prüfung bindet die genaue Fassung ihrer Gegenstände.
Ein gültiges Bildreview für einen alten Asset-Hash bestätigt
nicht das neue Bild. Eine Änderung von Zieltext oder
Voraussetzungen kann bisherige inhaltliche Evidenz berühren;
ein rein technischer Umzug kann unter einem eng begründeten
Kompatibilitätsnachweis anders zu behandeln sein. Beides darf
nicht durch pauschales Erneuern von Hashes verborgen werden.

\subsection{Veröffentlichung und Austausch}

Fachliche Änderungen erfolgen in kanonischen Autorenquellen.
Pakete, Bücher, Bilder und Validierungsberichte werden daraus neu gebaut;
historische veröffentlichte Archive bleiben unverändert.

JSON- und Ontologiepakete sollen aus demselben fachlichen Stand entstehen.
Das JSON-Laufzeitmodell ist die normative Referenz. Der Ontologie-Roundtrip
muss Zielidentitäten, Relationen, Rollen, Texte und weitere fachliche Felder
gemäß ihrem Gleichheitsvertrag erhalten. Unterschiedliche JSON-/RDF-Texte
benötigen semantische Gleichheit, binäre Ressourcen identische Bytes.
Das Austauschformat darf keine abweichende Laufzeitsemantik einführen.

Externe Zielreferenzen müssen vollständig aufgelöst sein. Ein Physikpaket
mit mathematischen Voraussetzungen benötigt deren Definitionen;
leere Platzhalter genügen nicht. Diese vollständige fachliche Closure
bestimmt noch nicht, welche Ziele für die Person als \code{target}
auswählbar sind.

Veröffentlichungs- und Prüfaussagen nennen Artefakt, Umfang, Version und
gebundene Quellen. Ein lokaler Build belegt weder fremde Hostanzeige noch
pädagogische Wirksamkeit. Dieses Kompendium erteilt keine Release- oder
Produktionsfreigabe.

\subsection{Lizenzen und private Lernergebnisse}

Software, technische Schemas, Entwicklerdokumentation und funktionale
Coachanweisungen stehen unter Apache-2.0; eigene Skill-Landschaften,
Zieltexte, didaktische Aufgaben, Karten, Bilder und Bildungsdokumentation
unter CC-BY-4.0. Entscheidend ist die Funktion, nicht die Dateiendung.
Gemischte Sammlungen sind nicht pauschal doppelt lizenziert.

Fremde Quellen, Abhängigkeiten, Medien und Marken behalten ihre Bedingungen.
Quellenangabe oder KI-Herkunft sind weder Rechtefreigabe noch Qualitätsbeweis.
Exporte führen zutreffende Lizenzen, Herkunftsangaben und Änderungshinweise
mit. Private Lernergebnisse, Sitzungsdaten und Gespräche gehören weder zur
öffentlichen Inhaltslizenz noch in öffentliche Quellen- oder Reviewpakete.

\paragraph{Quellen.}
\srcref{docs/concept/skill-graph/source-and-resource-links.md}{Quellen und Materialien};
\srcref{docs/concept/skill-graph/human-readable-source-rationales.md}{Menschenlesbare Quellenbegründungen};
\srcref{docs/qa-ci/graph-validation-rules.md}{Algorithmische Graphvalidierung};
\srcref{docs/qa-ci/curriculum-quality-maturity-and-routes.md}{Reifegrade und Qualitätsregeln};
\srcref{docs/concept/skill-graph/learning-goal-book-and-evidence-review-pipeline.md}{Lernzielbuch und Evidenzreview};
\srcref{docs/concept/skill-graph/dual-curriculum-package-releases.md}{Curriculum-Pakete};
\srcref{LICENSING.md}{Lizenzaufteilung}.

\clearpage
\ohead{Curriculare Planung}
% SPDX-License-Identifier: Apache-2.0
\section{Curriculare Level: Erwerbsordnung und Bewältigbarkeit}
\label{sec:levels}
\begin{principle}{Status: optionales Planungsprofil}
Die folgenden Bedingungen definieren einen curricularen Levelplan.
Sie ändern weder Grundgraphgültigkeit noch aktuelle Frontier oder
gespeicherte Mastery. Es wird damit kein bereits implementierter
Levelvalidator behauptet. Personalisierungsstufen, Schwierigkeit und
QA-Reifegrade sind andere Begriffe.
\end{principle}

Ein curricularer Level ist eine kontextgebundene Erwerbsetappe, etwa ein
Schuljahr, ein Semester oder eine fachübergreifende Stufe. Das anschauliche
Bild eines Wissensgebäudes hilft: Kompetenzen sind Bauteile, Voraussetzungen
tragende Beziehungen, ein Level eine geplante Bauetappe. Das Bild darf aber
nicht suggerieren, alle Fächer müssten vollständig abgeschlossen sein,
bevor irgendwo ein Ziel des nächsten Levels verfügbar wird.

\subsection{Plan, Eintrittsbasis und eindeutiger Ersterwerb}
Für einen aufgelösten Plan $\pi$ seien
\begin{align}
 T_\pi&\subseteq A &&\text{geplante Erwerbsziele},\\
 B_\pi&\subseteq A\setminus T_\pi &&\text{ausdrücklich angenommene Eintrittsbasis},\\
 n_\pi&\in\N_{>0},\qquad
 \ell_\pi:T_\pi\to\{1,\ldots,n_\pi\} &&\text{Erwerbslevel}.
\end{align}
Die Levelmengen und ihre Präfixe sind
\[
 L_i=\{g\in T_\pi\mid\ell_\pi(g)=i\},\quad
 L_{<i}=\bigcup_{j<i}L_j,\quad L_{\leq i}=\bigcup_{j\leq i}L_j.
\]
Jedes Ziel erhält in diesem Plan genau einen Ersterwerbslevel. Ein anderer
Bildungsgang darf es anders platzieren. Spätere Wiederholung, Vertiefung oder
Anwendung sind zusätzliche Aktivitäten, keine neue Zielidentität.
$B_\pi$ ist eine Planungsannahme; sie schreibt keine Lernergebnisse.
Persönliche Lücken in dieser Basis brauchen eine eigene Diagnose und
zusätzliche Lernzeit.

\begin{figure}[htbp]
\centering
\begin{tikzpicture}[node distance=10mm]
\node[goal,text width=24mm] (b) {Eintrittsbasis\\$B_\pi$};
\node[goal,text width=24mm,right=12mm of b] (l1) {Level 1\\$L_1$};
\node[goal,text width=24mm,right=12mm of l1] (l2) {Level 2\\$L_2$};
\node[goal,text width=24mm,right=12mm of l2] (l3) {Level 3\\$L_3$};
\draw[-{Stealth},draw=muted] (b)--(l1);
\draw[-{Stealth},draw=muted] (l1)--(l2);
\draw[-{Stealth},draw=muted] (l2)--(l3);
\coordinate (levelscenter) at ($(b.south)!0.5!(l3.south)$);
\node[info,anchor=north,text width=130mm] at ($(levelscenter)+(0,-7mm)$)
 {Jedes Level besitzt einen global gültigen Erwerbsweg\\und ein begründetes Zeitbudget.};
\end{tikzpicture}
\caption[Curriculare Erwerbsstufen]{Planungsreihenfolge, keine \code{requires}-Kanten zwischen
Levelcontainern. Einzelziele werden über ihre eigenen Voraussetzungen verfügbar.}
\label{fig:levels}
\end{figure}

\subsection{Strukturelle Machbarkeit durch einen Abschlusszeugen}
Ein Level $L_i$ ist strukturell machbar, wenn sich seine Ziele als
$a_1,\ldots,a_m$ aufzählen lassen, sodass
\begin{equation}
 a_j\in\Frontier\!\left(B_\pi\cup L_{<i}\cup
                         \{a_1,\ldots,a_{j-1}\}\right)
 \qquad(1\leq j\leq m).
 \label{eq:level-witness}
\end{equation}
Hier gilt die \emph{globale} Frontier. Vererbte und entlang offener Ziele verfolgte,
clusterbezogene und nach Referenzauflösung fachübergreifende Anforderungen
bleiben erhalten. Ein Normalmodusfilter würde gerade jene Anforderungen
ausblenden können, deren Einhaltung die Levelprüfung nachweisen soll.

Zur Diagnose dienen die benötigten Atome
\begin{equation}
 \ReqAtoms(g)=\bigcup_{\substack{p\in G\\(g,p)\in\Reff^+}}\Atoms(p).
 \label{eq:reqatoms}
\end{equation}
Aus einem Abschlusszeugen folgt notwendig
\begin{equation}
 \forall g\in T_\pi:\quad
 \ReqAtoms(g)\subseteq B_\pi\cup L_{\leq\ell_\pi(g)}.
 \label{eq:level-closure}
\end{equation}
Eine benötigte Kompetenz außerhalb $B_\pi\cup T_\pi$ ist ungeplant;
ein benötigtes Ziel mit höherem Erwerbslevel ist ein Reihenfolgefehler.
Voraussetzungen innerhalb desselben Levels sind erlaubt. Die Mengenbedingung
allein genügt bei Clusteranforderungen noch nicht: Der Selbstblocker aus
Abschnitt~\ref{sec:validity} erfüllt sie für $L_1=\{a\}$,
besitzt aber keinen Abschlusszeugen.

Ein endlicher Prüfalgorithmus startet für jedes Level mit
$M^{(0)}=B_\pi\cup L_{<i}$ und wiederholt
\begin{equation}
 M^{(r+1)}=M^{(r)}\cup\bigl(L_i\cap\Frontier(M^{(r)})\bigr).
\end{equation}
Er ist erfolgreich genau dann, wenn sein Fixpunkt ganz $L_i$ enthält.
Die Verfügbarkeit eines noch offenen Ziels kann durch zusätzliche
Erfüllungen nicht verschlechtert werden. Deshalb kann jede gleichzeitig
hinzugefügte Menge auch nacheinander erworben werden. Ein vorhandener
Abschlusszeuge kann umgekehrt nicht vorzeitig an diesem Verfahren scheitern.
Das Verfahren ist ein \emph{hypothetischer} Lerndurchlauf; es ändert keine
persönlichen Daten. Im rein atomaren DAG liefert eine voraussetzungsgetreue
topologische Ordnung den bekannten Spezialfall.

\subsection{Aufwand, Reserve und Budget}
Strukturelle Machbarkeit beweist noch nicht, dass ein Schuljahr zeitlich
ausreicht. Jede Zeitbehauptung nennt Referenzlernende, Eintrittsbedingungen,
Ziel- und Wahlauswahl, Kalender, Schätzmethode und Unsicherheit.
Alle Größen müssen in derselben kalibrierten Zeiteinheit vorliegen:
\begin{equation}
 \Work_\pi(i;\Theta_\pi)+\Reserve_\pi(i)\leq \Budget_\pi(i),
 \qquad i=1,\ldots,n_\pi.
 \label{eq:budget}
\end{equation}
Dabei sind $\Work\geq0$, $\Reserve\geq0$ und $\Budget>0$.
Reserve ist geschützt und wird nicht gleichzeitig als verplanbare Zeit
ausgegeben. Ein Nullpuffer ist möglich, aber offenzulegen. Bereits im
Arbeitsaufwand enthaltene Prüfungsvorbereitung oder Wiederholung darf nicht
ein zweites Mal als Reservebedarf gezählt werden.

Ein erklärtes Aktivitätenmodell kann
$\Work_\pi(i;\Theta_\pi)=\sum_{h\in\mathcal H_i}\widehat t_h(\Theta_\pi)$
verwenden. $\mathcal H_i$ enthält jeden tatsächlichen Aufwand genau einmal:
Ersterwerb, gezielte Übung, Abruf, Prüfung und weitere geplante Arbeit.
Dasselbe kanonische Atom wird nicht in Mathematik und Physik zweimal als
Ersterwerb gezählt; eine zusätzliche Anwendung in Physik kann dennoch
eigenständige Zeit kosten. Unkalibrierte \glqq Zielpunkte\grqq{} oder
\code{weight}-Werte sind keine Minutenschätzung.
Tabelle~\ref{tab:workload} veranschaulicht die getrennte Verbuchung von Arbeit und Reserve.

\begin{table}[htbp]
\centering
\begin{tabular}{@{}lrr@{}}
\toprule
Illustrative Planung & Stunden & Bedeutung\\
\midrule
Ersterwerb & 42 & Bereits um gemeinsame Ziele bereinigt\\
Übung und Anwendung & 12 & Zusätzliche Tätigkeiten\\
Prüfung und Wiederholung & 6 & Im Arbeitsaufwand enthalten\\
Geschützte Reserve & 10 & Nicht für obige Arbeit verplant\\
Verfügbare Zeit & 75 & Nach Kalenderausschlüssen\\
\bottomrule
\end{tabular}
\caption{Rechenbeispiel: $\Work=60$ und $\Reserve=10$ passen in
$\Budget=75$. Die erfundenen Zahlen demonstrieren nur die Rechnung,
keine empirisch abgesicherte Dauer eines SkillPilot-Curriculums.}
\label{tab:workload}
\end{table}

Ein Gesamtbudget genügt nicht, wenn Fächer, Räume oder Termine eigene
bindende Grenzen haben. Für eine zusätzlich deklarierte Fächerzuordnung
müssen auch die jeweiligen $\Work_{i,s}+\Reserve_{i,s}\leq \Budget_{i,s}$
gelten, ohne Doppelzählung gemeinsam zugeordneter Zeit. Ein Kalenderplan
muss außerdem Voraussetzungen vor der konkreten Anwendung bereitstellen.

\subsection{Gültigkeitsanspruch und Abschlussgrenzen}
Ein Levelplan ist unter diesem optionalen Profil gültig, wenn er sowohl
strukturell als auch unter seinen erklärten Annahmen zeitlich machbar ist.
Ohne belastbare Aufwandsschätzung ist allein die strukturelle Aussage
zulässig. Der Anspruch bindet Plan, Curriculumfassung, aufgelösten Umfang
und Schätzannahmen. Er garantiert keine identische Lerngeschwindigkeit
für jede Person.

Der Abschluss eines Levels wird aus seiner vollständigen Zielmenge und den
Zieltypverträgen abgeleitet. Er ist keine pauschale Voraussetzung für alle
Ziele des nächsten Levels. Soll ein Teil der Arbeit durch Prüfungsnachweise
ersetzt werden, benötigt das den folgenden gesonderten Abschlussvertrag.

\paragraph{Quellen.}
\srcref{docs/concept/skill-graph/curriculum-levels.md}{Curriculum Levels};
\srcref{docs/concept/skill-graph/graph-definition.md}{Graph Definition, Abschnitt 12};
\srcref{docs/concept/didactic/curriculum-time-axis-and-pacing.md}{Zeitachse und Kapazität}.

\clearpage
\ohead{Implizite Erledigung und Prüfungszugang}
% SPDX-License-Identifier: Apache-2.0
\section{Implizite Erledigung und direkter Prüfungszugang}
\label{sec:assessment}
\begin{principle}{Verbindliche Produktentscheidung vom 31. Oktober 2026}
Jedes gemeisterte Lernziel beendet seinen eigenen Vorbereitungsbedarf.
Ein von keinem offenen Ziel mehr benötigtes Nicht-Prüfungsziel ist implizit
erledigt und beendet ebenfalls seinen Bedarf. Jede Abschlussprüfung bleibt
ohne Mastery offen und muss direkt ansteuerbar sein. Manuelle
Freistellungsprofile und zusätzliche Prüfbelege sind dafür nicht erforderlich.
\end{principle}

Mastery bleibt der einzige gespeicherte Ergebniszustand pro atomarem Lernziel.
Implizite Erledigung wird abgeleitet und niemals als Mastery zurückgeschrieben.
Grundlage ist die \srcref{docs/concept/skill-graph/assessment-gated-completion.md}{aktuelle
Konzeptspezifikation}; der \srcref{docs/dev/assessment-gated-completion-pilot.md}{Implementierungsstand}
trennt lokale Tests, Deployment und tatsächliche Coach-Abnahme.

\subsection{Fester Umfang und direkte Bedarfsrelation}
Für eine Curriculumgeneration und das verbindliche Personal Curriculum sei
$T$ die Menge der eindeutigen atomaren Ziele des gewählten Umfangs.
Maßgeblich ist die aufgelöste Program-/Kompositionsstruktur. Ein temporärer
Lernfokus oder eingeklappter Baumzweig verändert diesen Umfang nicht.
Die Standardauswertung umfasst das Personal Curriculum; ein ausdrücklich
gewählter Abschnitt erhält seine eigene vollständige Auswertung.

Cluster-Voraussetzungen werden in atomare Pflichten expandiert;
vererbte Anforderungen und autoritative atomare \code{goalEntry}-Projektionen
werden berücksichtigt. Die vollständige Voraussetzungshülle $V_H$ enthält
auch versteckte und fachübergreifende Unterstützung. Unterstützende Ziele
werden dadurch nicht zu Zielen des Abschlussnenners.

Die Relation
\begin{equation}
 F_H=\{(u,v)\mid u\text{ benötigt unmittelbar }v
       \text{ nach atomarer Normalisierung}\}
\end{equation}
zeigt vom Verbraucher zur Voraussetzung. Fehlende Referenzen, Identitätsfehler
sowie Zyklen in \code{contains}, \code{requires} oder der expandierten
atomaren Relation sind Berechnungsfehler. Sie bedeuten nicht, dass keine
Arbeit mehr erforderlich sei.

\textbf{Eine statische transitive Hülle ist hier falsch.}
Bei $C\to B\to A$ beendet Mastery von B den Bedarf an A.
Eine zusätzliche direkte Kante $C\to A$ wäre dagegen eine eigenständige
Anforderung von C. Transitive Hülle und transitive Reduktion sind daher
keine bedeutungserhaltenden Umformungen dieser Abschlusssemantik.

Sei $Q_H\subseteq V_H$ die autoritativ bestimmte Menge der Abschlussprüfungen.
Die aktuelle Abbildung nutzt die bestehenden Prüfungs-/Assessmenttypen,
nicht Titel oder \code{coveredGoalIds}. Bereits im Kontext enthaltene
Prüfungspflichten verschwinden nicht dadurch, dass sie nur als Unterstützung
oder in einer anderen Ansicht sichtbar sind. Alternative Aufgabenformen
und Wahlentscheidungen müssen bereits bei der Umfangsauflösung geklärt sein.

\subsection{Mastery, Bedarf und kleinster Fixpunkt}
Mit der aktuellen globalen Mastery-Menge $M$ sei $E=M\cap V_H$.
Für eine Menge $X$ bereits ausdrücklich oder implizit erledigter Ziele gilt
\begin{align}
 \Consumers_H(v)&=\{u\mid(u,v)\in F_H\},\\
 \Active_H(X)&=\{(u,v)\in F_H\mid u\notin X\}.
\end{align}
Mastery beendet die \emph{eigenen ausgehenden} Anforderungen. Ein eingehender
Bedarf an einem bereits gemeisterten Ziel wird nicht durch dieses Ziel
hindurch an dessen frühere Vorbereitung weitergereicht.

Die drei Produktregeln ergeben unmittelbar
\begin{align}
 X_0&=E,\\
 X_{n+1}&=X_n\cup\{v\in V_H\setminus Q_H\mid
              \Consumers_H(v)\subseteq X_n\},\label{eq:waiver-step}\\
 X_*&=\bigcup_{n\geq0}X_n,\qquad I_H=X_*\setminus E.
 \label{eq:waiver-result}
\end{align}
Der Operator ist monoton und inflationär. Auf dem endlichen Graphen erreicht
die Iteration ihren kleinsten Fixpunkt. Dabei gilt stets
$I_H\cap E=\varnothing$ und $I_H\cap Q_H=\varnothing$.
Es entsteht keine zusätzliche gespeicherte Ergebnisart.

Die leere Verbrauchermenge ist kein Sonderfall: Ein gewöhnliches Endziel
ohne Verbraucher ist implizit erledigt. Dies gilt auch für Orientierung
und Gedächtnisziele. Die ausdrücklichen Interaktions- und Recall-Verträge
bleiben unverändert; weder ein Kartenerfolg noch eine motivational erfolgte
Interaktion wird durch die Auswertung erfunden. Enthält ein gültiger endlicher
Umfang keine Abschlussprüfung, wird seine gesamte nicht gemeisterte Arbeit
implizit erledigt. Zusätzliche Ausnahmen für einzelne Zielarten sind in
diesen drei Regeln nicht vorgesehen.

\paragraph{Äquivalenter linearer Algorithmus.}
Starte bei allen noch nicht gemeisterten Abschlussprüfungen. Folge ihren direkten
Voraussetzungskanten, besuche jedes Ziel höchstens einmal und stoppe an
\emph{jedem} gemeisterten Ziel. Die erreichten nicht gemeisterten Ziele
bleiben benötigt. Alle übrigen nicht gemeisterten Ziele sind implizit erledigt.
Auf einem DAG entspricht dies genau dem Fixpunkt und benötigt
$O(|V_H|+|F_H|)$ Zeit und Speicher. Eine quadratische Speicherung aller
indirekten Beziehungen entfällt.

\subsection{Beispiel: Mastery unterbricht den Bedarf}
\begin{figure}[htbp]
\centering
\begin{tikzpicture}[node distance=12mm]
\node[goal,fill=orange!12] (c) at (0,0) {C: Prüfung\\offen};
\node[goal,fill=green!9] (b) at (4.5,0) {B: Mastery};
\node[goal] (a) at (9,0) {A: implizit};
\draw[req,draw=warm,dashed] (c)--node[above,info]{benötigt B}(b);
\draw[req] (b)--node[above,info]{Bedarf endet}(a);
\node[info,text width=115mm] at (4.5,-1.4)
 {C bleibt offen. B ist schon gemeistert und benötigt A nicht mehr.\\
 Ohne weiteren offenen Verbraucher ist A implizit erledigt.};
\end{tikzpicture}
\caption{Mastery beendet Vorbereitung unabhängig von einem Prüfungserfolg.
Eine zusätzliche direkte Kante von C nach A würde A weiterhin erforderlich machen.}
\label{fig:waivers}
\end{figure}
\FloatBarrier

\begin{table}[htbp]
\centering
\begin{tabularx}{\linewidth}{@{}lX@{}}
\toprule
Fall & Ergebnis\\
\midrule
$C\to B\to A$, nur B gemeistert & A implizit; Prüfung C bleibt offen.\\
Zusätzlich $C\to A$ & A bleibt direkt für C erforderlich.\\
Prüfungen B und C benötigen A & Ein einzelner Erfolg reicht nicht, solange der andere Verbraucher offen ist.\\
Beide Prüfungen gemeistert & A und anschließend seine nicht mehr benötigte Vorbereitung werden implizit erledigt.\\
Gewöhnliches Endziel ohne Verbraucher & Implizit erledigt, auch ohne bisherigen Erfolg.\\
Prüfung ohne Verbraucher & Bleibt bis zu ihrer Mastery offen.\\
\bottomrule
\end{tabularx}
\caption{Bedarfsrechnung ohne manuelle Freistellungsregeln.}
\label{tab:waivers}
\end{table}

\subsection{Abschluss und direkte Prüfungsauswahl}
Der Umfang ist abgeschlossen genau dann, wenn
\begin{equation}
 \Complete(H)\iff T\subseteq X_*\ \land\ Q_H\subseteq E.
\end{equation}
Auf dem endlichen DAG sind alle Prüfungspflichten gemeistert auch hinreichend:
Ohne offene Prüfung bleibt keine dauerhafte Nachfrage nach Nicht-Prüfungsarbeit.
Der Nenner zählt trotzdem nur Ziele aus $T$, nicht zusätzliche Unterstützung.
Eine in Sekundarstufe~I implizit erledigte Kompetenz kann im größeren
stufenübergreifenden Umfang noch benötigt werden; ihre globale Mastery bleibt gleich.

Jede passende Abschlussprüfung ist ohne vorherige Voraussetzungsmastery direkt
wählbar. Bewusste Auswahl, Authentifizierung, Lernendenrevision,
Sitzungszuordnung und verfügbare Aufgaben bleiben erforderlich. Ein Fehlversuch
erzeugt keinen negativen Ergebnisdatensatz und nimmt einen früheren Erfolg
nicht zurück. Ausführliche KI-Bewertung mit Punkten, Begründungen und
Verbesserungshinweisen bleibt im Gespräch erhalten.

Die lesende Übersicht schreibt keine Navigation. Suche und seitenweise
Darstellung erschließen das gesamte Angebot; eine Seitengrenze darf keine
Prüfung dauerhaft unzugänglich machen. Empfehlungen dürfen den aktuellen
Abschnitt bevorzugen, ohne die verbindlichen Pflichten zu verkleinern.
Normale Pflichtvorschläge lassen implizit erledigte Arbeit aus. Freiwilliges
Lernen bleibt über die unveränderte fachliche Zugangsprüfung möglich.

\subsection{Cockpit, Lernplan und technische Grenzen}
Der Baum verwendet ein kompaktes Icon mit dem Tooltip
„Für diesen Abschnitt nicht mehr erforderlich“. Diese Aussage behauptet
keinen Prüfungserfolg, wenn ein gewöhnliches Lernziel den Bedarf beendet hat.
Mastery-Anzeigen und abgeleitete Erledigung bleiben unterscheidbar.

Für die feste, eindeutige Planbasis $G$ gilt standardmäßig
\begin{equation}
 I_{\mathrm{plan}}=|G\cap(M\cup I_H)|.
\end{equation}
Die automatische Bedarfsrechnung verwendet das vollständige Personal Curriculum;
erst anschließend wird mit der Planmenge $G$ geschnitten. Ein datierter Plan
ist keine eigene abgeschlossene Prüfungswelt: Ein reiner Vorbereitungsplan
darf nicht schon deshalb fertig sein, weil seine Prüfung außerhalb von $G$
liegt. Auch echte fachübergreifende Anforderungen bleiben Bedarf, während
Quoten und Fachbilanzen getrennt bleiben. Eine kleine Abschnittsauswertung
ändert die Plananrechnung nicht.
Die Cockpit-Einstellung „Nur explizit erreichte Lernziele zählen“ ist
standardmäßig aus; bei Auswahl zählt nur $|G\cap M|$.
Planbasis, Termine und Sollmengen ändern sich dadurch nicht. Der Periodenwert
$H$ zählt weiterhin ausschließlich tatsächliche, zeitlich belegte Abschlüsse;
aktuelle implizite Erledigung erzeugt keine historischen Lernerfolgsereignisse.

„Prüfungsabkürzungen anbieten“ steuert direkte Einstiege. Die Einstellung
verändert weder Mastery noch den berechneten Bedarf oder gültige Plananrechnung.
Der \srcref{docs/concept/didactic/unified-learning-plan-status.md}{gemeinsame
Lernplanstatusvertrag} gilt für Cockpit und Coach gleichermaßen.

Statische Kompilierungen werden begrenzt nach Curriculumgeneration,
Projektion, Zielumfang und Semantikrevision zwischengespeichert. Mastery und
implizite Ergebnisse sind keine zweite Wahrheitsquelle. Nach Ergebnis- oder
Umfangsänderungen wird neu gerechnet; alte implizite Mengen werden niemals
zu Mastery-Eingaben. Ungültige Daten bleiben als Einschränkung sichtbar.
Automatische Tests prüfen diese Verträge, ersetzen aber weder Aufgabenqualität
noch menschliche Erprobung oder nachgewiesene Lernwirksamkeit.

\clearpage
\ohead{Anhänge}
% SPDX-License-Identifier: Apache-2.0
\appendix
\section{Statusabgleich und Grenzen der Aussagen}
\label{sec:implementation}
Diese Fassung berücksichtigt die Produktentscheidung und den lokalen
Implementierungsstand vom 31. Oktober 2026. Die Aussagen über laufenden Code
sind keine Behauptung über eine verifizierte Produktivinstallation oder eine
vollständige Konformitätsprüfung. Frühere Dokumentationsbefunde werden nur
fortgeführt, soweit sie weiterhin zutreffen.

\begin{longtable}{@{}P{32mm}P{116mm}@{}}
\toprule
Gegenstand & Geprüfter Unterschied und Konsequenz\\
\midrule\endhead
Clustererfüllung & Die normative Formel berücksichtigt alle atomaren
Nachfahren. \code{LearnerService.computeEffectivePrereqMastery} und
\code{computeScopedPrerequisiteMastery} bevorzugen dagegen rekursiv
Kinder, die \code{isCoreForPrereqs} als Kern einstuft, sobald ein
solches Kind existiert. Diese Einstufung verwendet dort fehlende Tags
beziehungsweise \code{GK}-Tags, nicht einfach das Feld \code{core}.
Beide Regeln sind nicht identisch.\\
Binärer Zielzustand & Verbindlich ist geschafft oder noch nicht geschafft.
Neue allgemeine Schreibzugriffe und gewöhnliche Zielkontrollen verwenden
binäre Ergebnisse. Historische numerische Daten bleiben über die gemeinsame
Schwelle $0{,}9$ lesbar; daraus entsteht keine Teil-Mastery.\\
Gedächtnisabschluss & Der dokumentierte Zieltypvertrag fordert
verifizierten Abruf und keine fälligen Karten. Backend und
\code{useSrsMastery} verwenden derzeit allein den verifizierten Abruf
aller erforderlichen Karten. Erneute Fälligkeit und erneute
Fehlprüfung haben daher unterschiedliche Auswirkungen.\\
Graphvalidierung & Lokale direkte Doppelungen bleiben redundant.
Die frühere transitive Minimalitätsforderung ist aufgehoben: eine direkte
Kante kann trotz eines alternativen Pfades einen selbständigen Bedarf
tragen. Strukturtests bleiben von fachlicher Qualität zu unterscheiden.\\
Curriculare Level & Die mathematischen Machbarkeitsbedingungen sind ein
optionales Konzeptprofil. Es wird kein produktiv aktivierter neuer
Levelresolver oder Budgetvalidator behauptet.\\
Prüfungsabschluss & Automatische Bedarfsrechnung, allgemeiner direkter
Prüfungszugang und inklusive Plananrechnung sind lokal implementiert.
Die drei Produktregeln ersetzen die manuelle Profilpflege. Terminale
Test-, CI- und Deploymentnachweise sind getrennt zu prüfen.\\
\bottomrule
\end{longtable}

Die verbleibenden Unterschiede sind als Modell-Implementierungsfragen
kenntlich. Die bereits angepassten Laufzeitpfade dürfen nicht weiter als
unveränderte historische Abweichungen beschrieben werden.

\begin{center}\begin{minipage}{\linewidth}
\paragraph{Prüfbare Codeanker.}
\srcref{backend/src/main/java/com/skillpilot/backend/service/LearnerService.java}{LearnerService.java}:
\begin{itemize}\small
\item \code{isSrsMasteredByVerifiedRecall}
\item \code{computeEffectivePrereqMastery}
\item \code{computeScopedPrerequisiteMastery}
\item \code{isCoreForPrereqs}
\end{itemize}
\srcref{app/src/hooks/useSrsMastery.ts}{useSrsMastery.ts};
\srcref{app/scripts/validateGraph.ts}{validateGraph.ts};
\srcref{backend/src/main/java/com/skillpilot/backend/domain/Mastery.java}{Mastery.java};
\srcref{app/src/components/GoalCard.tsx}{GoalCard.tsx};
\srcref{docs/qa-ci/graph-validation-rules.md}{Graph Validation Rules}.
Die Links zeigen auf die weiterentwickelte Hauptlinie; für eine historische
Prüfung ist zusätzlich der genannte Stand zu verwenden.
\end{minipage}\end{center}

\section{Prüfstrategie für die Abschlussberechnung}
Das begleitende \code{latex/check-assessment-model.py} ist eine
ausführbare mathematische Referenzprüfung mit ausschließlich
Python-Standardbibliothek. Es liest keine Lernerdaten und verändert
keinen Runtimezustand. Es ergänzt die allgemeinen Beweise durch
systematische endliche Gegenbeispielsuche.

\begin{center}\begin{tabularx}{\linewidth}{@{}P{40mm}X@{}}
\toprule
Prüfbereich & Geprüfte Aussage\\
\midrule
Konkrete Fälle & Mastery-Barriere, zusätzliche unabhängige direkte Kante,
geteilte Grundlagen, gewöhnliche Endziele und offene Prüfungen.\\
Mastery als einzige Eingabe & Keine impliziten Prüfungserfolge;
Mastery bleibt Teil der erledigten Menge. Mehr Mastery öffnet bei
unverändertem Kontext keine Arbeit erneut.\\
DAGs bis vier Knoten in topologischer Ordnung & Alle Kantenmengen und
jede Mastery-/Prüfungsbelegung: insgesamt 16\,932 Fälle. Gleichheit von
lokalem Verbraucherfixpunkt und Traversierung ab offenen Prüfungen.\\
Ungültige Eingaben & Zyklen und unbekannte Referenzen bleiben Fehler,
keine leere und damit vermeintlich abgeschlossene Arbeitsmenge.\\
\bottomrule
\end{tabularx}\end{center}

Diese Prüfung enthält keine Modelle realer Aufgabenqualität,
Benotungszuverlässigkeit, Identitätskontrolle oder pädagogischer Wirkung.
Runtime-Regressionen prüfen zusätzlich konkurrierende Änderungen,
idempotente Mastery-Speicherung, Rechte, korrigierte Mastery und
übereinstimmende Anzeigen prüfen. Der Referenzalgorithmus
ist kein neuer produktiver Service.

\section{Begriffe und Symbole}
\begin{longtable}{@{}P{36mm}P{112mm}@{}}
\toprule
Zeichen / Begriff & Bedeutung\\
\midrule\endhead
$G,A,K$ & Ziele, strukturelle Atome, Cluster.\\
$C,R,\Reff$ & Enthaltensein, direkte Voraussetzungen,
direkte plus geerbte Voraussetzungen.\\
$X^+,X^{-1}$ & Nichtleere transitive Hülle und Umkehrrelation.\\
$\Anc,\Desc$ & Vorfahren und Nachfahren; in der Quellspezifikation
\code{Ancestors} und \code{Descendants}.\\
$\Atoms(g)$ & Atomare Basis; bei einem Atom die einelementige Menge.\\
$O,Z,E_{\mathrm{route}}$ & Orientierungsanker, autonome Endziele und
explizite Routenausnahmen; $O,Z$ entsprechen dort $M,T$.\\
$R_A$ & Direkte atomare Voraussetzungskanten $R\cap(A\times A)$.\\
$M_A,m,\Sat$ & Geschaffte Atome, ihr binärer Indikator
$m(a)\in\{0,1\}$ und daraus abgeleitete Erfüllung eines Ziels.\\
$F,G_F$ & Filterprädikat und ausgewählter Zielbereich.\\
$D,V_d$ & Filterdimensionen und Werte einer Dimension $d$.\\
Frontier & Noch nicht erfüllte atomare Ziele mit erfüllten
Voraussetzungen im gewählten Prüfmodus.\\
Personal Curriculum & Verbindlicher längerfristiger persönlicher
Zielumfang; Personalisierungslevel 2.\\
Fokus & Vorübergehende Arbeitsauswahl im Personal Curriculum;
Personalisierungslevel 3.\\
$T_\pi,B_\pi,\ell_\pi$ & Zielmenge, Eintrittsbasis und Erwerbslevel
eines curricularen Plans.\\
$T,Q_H$ & Ziele des gewählten Umfangs und Abschlussprüfungen seines vollständigen Kontextes.\\
$F_H$ & Direkte atomar expandierte und geerbte Bedarfsrelation;
keine statische transitive Hülle.\\
$H=(V_H,F_H)$ & Vollständiger Berechnungskontext einschließlich unterstützender Voraussetzungen.\\
$E$ & Vorhandene Mastery im Kontext: $E=M\cap V_H$.\\
$X_*,I_H$ & Berechnete explizit oder implizit erledigte Arbeit;
$I_H=X_*\setminus E$ enthält ausschließlich implizite Erledigung.
Keine zusätzlichen gespeicherten Ergebniszustände.\\
M7 & Benannter maschineller Curriculum-QA-Reifegrad,
kein menschlicher Wirksamkeitsnachweis.\\
\bottomrule
\end{longtable}

\section{Quellenzuordnung und Abbildungen}
Die Kapitel sind eine redaktionell zusammengeführte deutsche Darstellung
der folgenden Repositoryquellen. Querverweise nennen die jeweilige
fachliche Verantwortung; sie sind keine Behauptung, dass alle optionalen
Konzepte bereits implementiert seien.

\begin{longtable}{@{}P{38mm}P{110mm}@{}}
\toprule
Inhalt & Maßgebliche Quelle\\
\midrule\endhead
Formaler Kern und Filter &
\srcref{docs/concept/skill-graph/graph-definition.md}{Graph Definition},
Abschnitte 1--11.\\
Schichten und Migration &
\srcref{docs/concept/skill-graph/general-goal-system-and-migration.md}{General Goal System and Migration};
\srcref{docs/concept/skill-graph/canonical-gymnasium-rollout.md}{Canonical Gymnasium Rollout}.\\
Sichten und persönliche Auswahl &
\srcref{docs/concept/skill-graph/view-projection-and-goal-placement.md}{View Projection and Goal Placement};
\srcref{docs/concept/levels-of-personalization.md}{Levels of Personalization}.\\
Zieltypen und Evidenz &
\srcref{docs/concept/skill-graph/node-types.md}{Node Types};
\srcref{AGENTS.md}{AGENTS.md, insbesondere Orientierung und Mastery}.\\
Curriculare Planung &
\srcref{docs/concept/skill-graph/graph-definition.md}{Graph Definition, Abschnitt 12};
\srcref{docs/concept/skill-graph/curriculum-levels.md}{Curriculum Levels};
\srcref{docs/concept/didactic/curriculum-time-axis-and-pacing.md}{Curriculum Time Axis and Pacing}.\\
Impliziter Abschluss &
\srcref{docs/concept/skill-graph/assessment-gated-completion.md}{Assessment-Gated Completion};
\srcref{docs/concept/runtime-workflows/exam-preparation-planning.md}{Exam Preparation Planning}.\\
Herkunft und Materialien &
\srcref{docs/concept/skill-graph/source-and-resource-links.md}{Source and Resource Links};
\srcref{docs/concept/skill-graph/content-integration.md}{Content Integration};
\srcref{docs/concept/skill-graph/human-readable-source-rationales.md}{Human-Readable Source Rationales}.\\
Bilder und Zielbücher &
\srcref{docs/concept/skill-graph/atomic-goal-visualizations.md}{Atomic Goal Visualizations};
\srcref{docs/concept/skill-graph/learning-goal-book-and-evidence-review-pipeline.md}{Learning-Goal Book and Evidence-Review Pipeline}.\\
Validierung und QA &
\srcref{docs/qa-ci/graph-validation-rules.md}{Graph Validation Rules};
\srcref{docs/qa-ci/curriculum-quality-maturity-and-routes.md}{Curriculum Quality Maturity and Routes};
\srcref{docs/qa-ci/semantic-atomicity-review.md}{Semantic Atomicity Review}.\\
Austausch und Rechte &
\srcref{docs/concept/skill-graph/dual-curriculum-package-releases.md}{Dual Curriculum Package Releases};
\srcref{docs/concept/skill-graph/mem-fwu-roundtrip-plan.md}{MEM/FWU Roundtrip Plan};
\srcref{LICENSING.md}{LICENSING.md}.\\
\bottomrule
\end{longtable}

Die sieben didaktischen Graphabbildungen stammen aus
\href{https://github.com/enpasos/skillpilot/tree/main/docs/concept/skill-graph/assets/graph-definition}{SkillPilot: Graph Definition, Abbildungsbestand}
und werden unverändert mit deutschen Erläuterungen verwendet.
Die Diagramme auf dem Titelblatt sowie zu Levelplanung und Freistellung
entstehen direkt in TikZ. Maßgeblich sind die Formeln und Bildlegenden;
die Bilder behaupten keine zusätzlichen Kanten oder Anforderungen.

Technischer Text, Formelspezifikation und Build-Werkzeuge stehen unter
Apache-2.0; die didaktischen Illustrationen unter CC BY 4.0, soweit die
entsprechenden Rechte vorliegen. Lizenztexte und Abgrenzung:
\srcref{LICENSE}{LICENSE},
\srcref{LICENSES/CC-BY-4.0.txt}{LICENSES/CC-BY-4.0.txt} und
\srcref{LICENSING.md}{LICENSING.md}.
Quellenhinweis für die Illustrationen: SkillPilot, enpasos -- Enterprise
Patterns \& Solutions GmbH, Graph Definition, Fassung vom 29. Oktober 2026.
Zusammenführung und deutsche Erläuterungen: diese Dokumentfassung.

\section{Dokument bauen und prüfen}
Die Hauptdatei \code{skillpilot-graph-definition.tex} bindet die Kapitel
aus \code{latex/} und die PNG-Abbildungen aus
\code{assets/graph-definition/} ein. Die neuen Diagramme sind TikZ-Vektorgrafiken.

Vom Repositorywurzelverzeichnis aus:
\begin{quote}\small\ttfamily
python3 docs/concept/skill-graph/latex/build-document.py\par
python3 docs/concept/skill-graph/latex/check-assessment-model.py
\end{quote}
Benötigt werden Python 3 und LuaLaTeX mit deutscher Sprachunterstützung
und den Präambelpaketen. Hilfsdateien liegen unter
\code{tmp/skill-graph-document/}, das PDF neben der Hauptdatei.
Der Build löst Verweise in mehreren Läufen auf und bindet Eingabedateien
und PDF durch SHA-256 im Build-Manifest samt Datum und Werkzeugversion.
Über unterschiedliche TeX-Versionen hinweg ist Byteidentität nicht garantiert.

Zum Abschluss gehören ein fehlerfreier LaTeX-Lauf, aufgelöste Verweise,
die Prüfung auf überbreite beziehungsweise überhohe Boxen und eine
visuelle Sichtung des erzeugten PDF. Die Formelprüfung ergänzt diese
Dokumentprüfungen, ersetzt aber weder Runtime-Tests noch unabhängige
fachliche und empirische Validierung.

\end{document}
